% =====================================================================
%  ANEXO C. INSTRUMENTOS DE RECOLECCIÓN DE DATOS Y FORMATOS DEL SGC
%  C.1 (lista de verificación ISO 9001:2026) se incluye desde
%  capitulos/91a_lista_verificacion_iso9001.tex (otro archivo).
%  C.2 a C.8: formatos listos para imprimir (tamaño carta).
% =====================================================================
\chapter{INSTRUMENTOS DE RECOLECCIÓN DE DATOS Y FORMATOS DEL SGC}
\label{anx:instrumentos}

% --- Utilidades locales para los formatos -----------------------------
\providecommand{\casilla}{\framebox[1.9ex]{\rule{0pt}{1.3ex}}}
\providecommand{\campo}[1]{#1~\hrulefill}
\providecommand{\fila}{\rule{0pt}{0.58cm}}
\providecommand{\firma}[1]{\parbox[b]{3.7cm}{\centering\rule{0pt}{0.75cm}\rule{3.4cm}{0.4pt}\\[-1pt]{\scriptsize #1}}}
\newcommand{\cabeceraformato}[2]{%
  \par\noindent{\footnotesize\setstretch{1}%
  \renewcommand{\arraystretch}{1.35}%
  \begin{tabular}{@{}|p{0.30\linewidth}|p{0.44\linewidth}|p{0.183\linewidth}|@{}}
  \hline
  \textbf{UATF -- Edificio de Posgrado} & \multirow{2}{=}{\centering\textbf{#2}} & \textbf{Código:} #1 \\
  Ítem 13: losa alivianada h~=~25~cm & & Versión 0 \\
  \hline
  \end{tabular}\par}\smallskip}

Este anexo reúne los instrumentos de recolección de datos de la investigación y los formatos del Sistema de Gestión de la Calidad (SGC) propuesto en la Sección~\ref{sec:diseno-sgc}. Los formatos C.2 a C.6 son los registros del SGC: se imprimen en tamaño carta, se llenan en obra y se archivan por losa y por paño. Los instrumentos C.7 y C.8 sirven a la investigación: la guía de entrevista para el diagnóstico (OE1) y el formato de validación por juicio de expertos (OE5). La relación de cada formato con el sistema es la siguiente:

\begin{center}
\footnotesize\setstretch{1}\renewcommand{\arraystretch}{1.25}
\begin{tabularx}{\linewidth}{@{}P{1.0cm}LP{4.6cm}P{3.0cm}@{}}
\toprule
\textbf{Código} & \textbf{Formato o instrumento} & \textbf{Componente del SGC} & \textbf{Responsable} \\
\midrule
C.1 & Lista de verificación de requisitos ISO~9001:2026 & Diagnóstico y auditoría interna & Investigador; auditor \\
C.2 & Fichas PPI de liberación por paño & Plan de puntos de inspección (Tabla~\ref{tab:ppi}) & Encargado de calidad; Supervisor \\
C.3 & Ficha de recepción de materiales & PCM-01 (Tabla~\ref{tab:pcm01}) & Almacenero; encargado de calidad \\
C.4 & Control de materiales y conciliación por paño & PCM-02, PCM-03 y PCM-04 & Almacenero; encargado de calidad \\
C.5 & Registro de no conformidades y acción correctiva & PCM-05; procedimiento de NC & Encargado de calidad \\
C.6 & Planilla de costos de la calidad (PAF) & Registro de costos de la calidad & Encargado de calidad \\
C.7 & Guía de entrevista semiestructurada & Diagnóstico (OE1) & Investigador \\
C.8 & Formato de validación por juicio de expertos & Validación (OE5) & Investigador; expertos \\
\bottomrule
\end{tabularx}
\end{center}

\IfFileExists{capitulos/91a_lista_verificacion_iso9001.tex}{% =====================================================================
%  ANEXO C.1. LISTA DE VERIFICACIÓN DE REQUISITOS DE ISO 9001:2026
%  APLICADA (DOCUMENTALMENTE) A LA EJECUCIÓN DE LA PRIMERA LOSA
%  Se inserta dentro del Anexo C (capitulos/91_anexo_instrumentos.tex)
%  con % =====================================================================
%  ANEXO C.1. LISTA DE VERIFICACIÓN DE REQUISITOS DE ISO 9001:2026
%  APLICADA (DOCUMENTALMENTE) A LA EJECUCIÓN DE LA PRIMERA LOSA
%  Se inserta dentro del Anexo C (capitulos/91_anexo_instrumentos.tex)
%  con % =====================================================================
%  ANEXO C.1. LISTA DE VERIFICACIÓN DE REQUISITOS DE ISO 9001:2026
%  APLICADA (DOCUMENTALMENTE) A LA EJECUCIÓN DE LA PRIMERA LOSA
%  Se inserta dentro del Anexo C (capitulos/91_anexo_instrumentos.tex)
%  con \input{capitulos/91a_lista_verificacion_iso9001}.
%  Las calificaciones se reproducen en el script
%  03_Datos_Obra_Edificio_Posgrado/04_Extraccion_de_Datos/scripts/
%  cap2_diagnostico_figuras.py (Figura de brechas del Capítulo 2).
% =====================================================================
\section*{C.1. Lista de verificación de requisitos de ISO~9001:2026 aplicada a la ejecución de la primera losa}
\addcontentsline{toc}{section}{C.1. Lista de verificación de requisitos de ISO~9001:2026 (primera losa)}
\label{anx:lista-iso}

\textbf{Propósito y alcance.} La lista sirve para diagnosticar las brechas entre la gestión real del proceso constructivo de la losa alivianada y los requisitos de las cláusulas 4 a 10 de la norma ISO~9001:2026 \parencite{iso9001_2026}. Se aplicó a la ejecución de la primera losa (41 paños; 1.109,84~\mdos{}), entendida como el sistema formado por la empresa contratista, su Dirección de Obra y los controles de la Supervisión que constan en los documentos. Los requisitos se redactaron como paráfrasis del autor, en términos del proceso de la losa; no reproducen el texto de la norma. La numeración corresponde a la estructura de cláusulas que la edición de 2026 conserva de la de 2015. \pendiente{contrastar la numeración de las subcláusulas con el texto de la NB/ISO~9001:2026 de IBNORCA \parencite{ibnorca9001_2026}.}

\textbf{Método.} La aplicación fue documental y la realizó el investigador. Se revisaron, página por página, la Planilla de avance de obra N.\textsuperscript{o}~2 con sus anexos (informe de Supervisión, cómputos, certificados, contrato, Libro de Órdenes y cronogramas), la planilla y los cómputos de cierre por resolución del contrato, el expediente del ítem~13, las especificaciones técnicas y los planos 01 a 03. Cada requisito aplicable se calificó así: \textbf{1} = cumple (el requisito está establecido, se aplicó y hay evidencia documental); \textbf{0,5} = cumple parcialmente (hay evidencia de una aplicación parcial, informal o sin registros completos); \textbf{0} = no cumple (no se encontró evidencia en los documentos revisados, o la evidencia muestra lo contrario); \textbf{NA} = no aplica al proceso de la losa (se excluye del denominador). El porcentaje de cada cláusula es la suma de las calificaciones dividida entre el número de requisitos aplicables. La interpretación usa la escala de \textcite[p.~42]{mamanicalsina2025}: 0--25~\% deficiente o no implementado; 26--50~\% básico o en desarrollo; 51--75~\% aceptable; 76--100~\% excelente.

\textbf{Limitación.} Una calificación de 0 por falta de evidencia documental no prueba que la práctica no existiera; prueba que no quedó registrada. Las entrevistas de la guía C.7 a la Dirección de Obra, a la Supervisión y a la Fiscalización permitirán verificar o corregir estas calificaciones. Para las auditorías internas del SGC propuesto, la columna de requisitos se usa sin cambios y las columnas de evidencia y calificación se llenan en cada aplicación.

\begingroup\footnotesize\setstretch{1}\setlength{\tabcolsep}{3pt}
\begin{longtable}{@{}P{1.15cm}P{3.85cm}P{5.05cm}Q{0.85cm}P{4.15cm}@{}}
\caption{Lista de verificación C.1 aplicada a la ejecución de la primera losa}\label{tab:lista-iso}\\
\toprule
\textbf{Cláusula} & \textbf{Requisito (paráfrasis aplicada a la losa)} & \textbf{Evidencia documental encontrada} & \textbf{Calif.} & \textbf{Observación} \\
\midrule
\endfirsthead
\multicolumn{5}{@{}l}{\textit{(continuación)}}\\
\toprule
\textbf{Cláusula} & \textbf{Requisito (paráfrasis aplicada a la losa)} & \textbf{Evidencia documental encontrada} & \textbf{Calif.} & \textbf{Observación} \\
\midrule
\endhead
\midrule
\multicolumn{5}{@{}r}{\textit{(continúa)}}\\
\endfoot
\bottomrule
\endlastfoot
%
\multicolumn{5}{@{}l}{\textbf{4. Contexto de la organización}}\\
4.1 & Determinar las cuestiones externas e internas que afectan al proceso de la losa (abastecimiento, precios, clima y pertinencia del cambio climático) y darles seguimiento. & No hay análisis de contexto. Las lluvias (siete días) y la demora de materiales por escasez de combustible solo se anotaron después de ocurridas, para pedir ampliación de plazo (PAO2, pp.~114--124). & 0 & La escasez de diésel ya era pública en junio de 2025 (FANCESA; Las Lomas) y no generó previsiones. \\
4.2 & Identificar las partes interesadas pertinentes (UATF, Supervisión, Fiscalización, proveedores, laboratorio, usuarios) y sus requisitos. & El contrato define partes y obligaciones (PAO2, pp.~91--112); la ET y los planos fijan los requisitos del producto. No hay análisis de las demás partes interesadas. & 0,5 & Documentados solo los requisitos de la entidad contratante. \\
4.3 & Definir y documentar el alcance del sistema, con los requisitos que no se aplican y su justificación. & No existe un sistema declarado ni un documento de alcance; los 15 documentos de la planilla no incluyen un plan de la calidad (PAO2, pp.~3--4, 12). & 0 & --- \\
4.4 & Determinar los procesos de la losa (abastecimiento, montaje, armado, hormigonado, curado y medición) con entradas, salidas, secuencia, criterios, responsables y recursos. & La secuencia general figura en la ET (pp.~1--4) y en el cronograma A-8 (PAO2, p.~137); no hay mapa de procesos, procedimientos ni responsables por actividad. & 0,5 & El cronograma no refleja la secuencia real: la losa se armó en noviembre de 2025. \\
%
\multicolumn{5}{@{}l}{\textbf{5. Liderazgo}}\\
5.1.1 a & Demostrar el compromiso de la dirección de la empresa con la eficacia del sistema (rendición de cuentas, recursos, integración en los procesos). & La dirección de la empresa solo aparece en las cartas de entrega de planillas y en las firmas de cobro (PAO2, pp.~12, 14, 88). & 0 & --- \\
5.1.1 b & Promover una cultura de la calidad y un comportamiento ético en la ejecución, la medición y el registro de la obra. & No hay evidencia de acciones; las diferencias de cantidades del acero no tienen explicación documentada (PAO3, hoja de cómputos, filas 375--541; Exp., p.~12). & 0 & Énfasis de la edición 2026. Es ausencia de evidencia, no un juicio sobre las personas. \\
5.1.2 & Mantener el enfoque al cliente: cumplir los requisitos de la UATF y los legales, y atender los riesgos que afectan la conformidad de la losa. & Obra ejecutada con autorización de la Supervisión (PAO2, pp.~114--125), pero con retraso acumulado de 3,03~\% (p.~10) y con garantías, ensayos y dosificaciones pendientes al cierre del periodo (p.~9). & 0,5 & --- \\
5.2 & Establecer, comunicar y mantener disponible una política de la calidad apropiada al proyecto. & No se encontró. & 0 & --- \\
5.3 & Asignar y comunicar roles, responsabilidades y autoridades, incluido quién informa sobre el desempeño del sistema. & Director de Obra designado y evaluado (PAO2, p.~9); Supervisora y Fiscal designados (pp.~7, 127). No hay organigrama ni responsable de calidad; falta la constancia del Director de Obra en las 12 hojas del Libro de Órdenes (pp.~114--125). & 0,5 & Cl.~25.ª del contrato (PAO2, p.~100). \\
%
\multicolumn{5}{@{}l}{\textbf{6. Planificación}}\\
6.1 a & Determinar los riesgos y las oportunidades del proceso (suministro, precios, clima, cambios de diseño). & No hay matriz de riesgos. Se materializaron riesgos no previstos: viguetas inexistentes en el mercado en los tamaños requeridos (PAO2, p.~11), entrega del 30~\% del acero pedido (carta de Las Lomas, Exp., p.~27) y demoras por combustible (PAO2, p.~123). & 0 & --- \\
6.1 b & Planificar acciones proporcionales a los riesgos, integrarlas en los procesos y evaluar su eficacia. & Sin evidencia; los eventos solo se tramitaron como causales de ampliación de plazo (PAO2, pp.~114--124), aunque la cl.~20.ª no los considera fuerza mayor (p.~96). & 0 & --- \\
6.2 & Establecer objetivos de la calidad medibles (desperdicio, conformidad del hormigón, plazos de liberación) y planificar cómo lograrlos. & El único objetivo medido es el avance físico-financiero (PAO2, p.~10). & 0 & --- \\
6.3 & Planificar y controlar los cambios que afectan al sistema o al producto. & Los cambios contractuales se tramitaron formalmente: contrato modificatorio N.\textsuperscript{o}~1 y dos ampliaciones de plazo (PAO2, pp.~4, 7). Los cambios técnicos de la losa (viguetas, complemento) no tienen registro (p.~11; p.~81; ET, p.~3). & 0,5 & --- \\
%
\multicolumn{5}{@{}l}{\textbf{7. Apoyo}}\\
7.1.2 y 7.1.3 & Proveer las personas y la infraestructura necesarias (frentes de trabajo, apuntalamiento, mezclado, transporte y vibrado del hormigón). & Dos frentes de obra civil y dos de vaciado de $\approx$~6 trabajadores (PAO2, p.~9); apuntalamiento de madera, carretillas y vibrador en las fotografías (pp.~74--76). La Supervisión pidió más cuadrillas para recuperar el retraso (p.~11). & 0,5 & Recursos presentes, pero insuficientes para el cronograma. \\
7.1.4 & Mantener un ambiente adecuado para la operación (protección del hormigón ante lluvia y bajas temperaturas). & No hay previsiones de vaciado ni de curado en tiempo frío o con lluvia; solo el registro de las lluvias (PAO2, pp.~114--124). & 0 & ACI~318-25 exige un plan de hormigonado en tiempo frío (p.~583). \\
7.1.5 & Disponer de recursos de seguimiento y medición adecuados y de un muestreo confiable (laboratorio, moldes, cono de Abrams, balanzas). & Ensayos en el laboratorio de la UATF (PAO2, p.~78), con probetas «traídas y elaboradas por la empresa solicitante»; no constan equipos de obra ni su verificación. & 0,5 & --- \\
7.1.6 & Determinar y mantener los conocimientos necesarios (instrucciones de montaje del fabricante, lecciones aprendidas). & El certificado de PRETENSA remite a los pasos de montaje del fabricante (PAO2, p.~80); no consta que estuvieran en obra. & 0 & --- \\
7.2 & Determinar la competencia necesaria (armadores, encofradores, personal de vaciado y muestreo), asegurarla y conservar evidencia. & Solo se evalúa al Director de Obra («Buena», PAO2, p.~9); la Supervisión trabajó con apoyo de pasantes (p.~10); no hay registros de competencia ni de capacitación. & 0,5 & --- \\
7.3 & Asegurar que el personal conozca la importancia de su trabajo para la calidad y las consecuencias de incumplir. & Sin evidencia; las fotografías muestran la parrilla pisada durante el vaciado y sin separadores visibles (PAO2, p.~76). & 0 & Observación fotográfica a confirmar. \\
7.4 & Definir la comunicación interna y externa (qué, cuándo, a quién, cómo y quién comunica). & El Libro de Órdenes funcionó como canal formal: 24 asientos con respuesta el mismo día (PAO2, pp.~114--125); hay cartas formales (pp.~2, 5, 12, 88). Encabezados vacíos y sin constancia del Director de Obra. & 0,5 & Sin evidencia de comunicación interna con el personal. \\
7.5.2 & Crear y actualizar la información documentada con identificación, fecha, revisión y aprobación. & Planillas, cómputos y croquis firmados y sellados (PAO2, pp.~14, 18--72), pero con errores: rótulo ZT10 repetido (p.~18), columna tipo~5 con 6 y 7 unidades (p.~38), hoja duplicada (pp.~51--52), anticipo de 10~\% en un contrato sin anticipo (p.~14), desfase del periodo mal calculado (p.~10). & 0,5 & --- \\
7.5.3 a & Controlar la información documentada: versión vigente, coherencia entre documentos y normas actualizadas. & Tres geometrías del complemento: 100~$\times$~50~$\times$~20~cm (ET, p.~3), altura de 0,15~m (Plano~02) y 20~$\times$~100~$\times$~44~cm (PAO2, p.~81); la ET cita la CBH-87 y NB antiguas (pp.~1--2); dos versiones de la planilla de cierre con distinto descuento de acero (Exp., p.~2; PAO3); el archivo de cierre conserva una hoja de otro contrato (PAO3). & 0 & --- \\
7.5.3 b & Conservar los registros que evidencian la conformidad (recepción, inspección, ensayos, consumos, vaciado y curado). & No se encontraron registros de recepción, inspección del armado, consumos, vaciado ni curado de la losa (PAO2, pp.~1--139; Exp., pp.~1--29). & 0 & Brecha central para la línea base. \\
%
\multicolumn{5}{@{}l}{\textbf{8. Operación}}\\
8.1 a & Planificar el proceso de la losa: criterios de aceptación, controles y registros necesarios (plan de la calidad, PPI). & La ET fija criterios: f'c~=~250~kg/cm\textsuperscript{2}, cemento $\geq$~400~kg/\mtres{}, vaciado por franjas y monolítico con las vigas (pp.~1--4). No hay plan de la calidad ni PPI. & 0,5 & --- \\
8.1 b & Determinar los materiales necesarios por paño (cantidades teóricas para pedido y control). & El cómputo del ítem~13 solo lista las áreas de los 41 paños (PAO2, pp.~28--29); no hay cantidades de viguetas, EPS, hormigón ni acero por paño. & 0 & --- \\
8.2.2 y 8.2.3 & Determinar y revisar los requisitos del producto (contrato, ET, planos) antes de ejecutar, y resolver las diferencias. & Requisitos disponibles (PAO2, pp.~91--112; ET; planos 01 a 03), pero las discrepancias del complemento, de su altura y de las normas citadas no se resolvieron, pese a la consulta previa que exige la cl.~33.1 (p.~107). & 0,5 & --- \\
8.2.4 & Actualizar la información y comunicarla cuando cambian los requisitos. & «Modificaciones estructurales» por falta de viguetas de los tamaños requeridos, sin documento que las describa (PAO2, p.~11); EPS distinto de la ET sin orden de cambio (p.~81). & 0 & --- \\
8.3 & Diseño y desarrollo de los productos. & --- & NA & El diseño estructural corresponde al proyectista de la UATF. \\
8.4.1 & Aplicar criterios para evaluar, seleccionar y hacer seguimiento a los proveedores externos. & Sin criterios documentados ni reevaluación tras los incumplimientos (acero entregado en 30~\%, Exp., p.~27; viguetas no disponibles, PAO2, p.~11). La certificación NB~997 e ISO~9001 de PRETENSA (p.~80) es un atributo del proveedor, no un criterio registrado. & 0 & --- \\
8.4.2 & Verificar lo suministrado antes de incorporarlo a la obra (certificados, inspección de recepción, ensayos). & Certificados de PRETENSA y UNICON del 27/11/2025, sin firma y posteriores al inicio del armado (14/11/2025) (PAO2, pp.~80--81, 122); certificados de acero de $\phi$20 a $\phi$8 (pp.~82--86), sin $\phi$6 ni $\phi$10; sin certificados de cemento ni áridos; sin inspección de recepción. & 0,5 & Cl.~26.6 (PAO2, p.~102); ACI~318-25 (p.~587). \\
8.4.3 & Comunicar a los proveedores los requisitos de lo que se compra (tipo, dimensiones, cantidad, aprobación). & No hay órdenes de compra; el EPS entregado (20~$\times$~100~$\times$~44~cm; 2.350 piezas) no coincide con la dimensión de la ET (PAO2, p.~81; ET, p.~3). & 0 & --- \\
8.5.1 a & Ejecutar en condiciones controladas: información disponible, seguimiento en las etapas apropiadas y personal competente. & Verificación del armado antes de cada vaciado, registrada en el Libro de Órdenes (PAO2, pp.~114--125), sin lista de verificación ni criterios; no consta la fecha del hormigonado de la losa (p.~125). & 0,5 & --- \\
8.5.1 b & Controlar el hormigón y validar el proceso de hormigonado: dosificación aprobada, asentamiento, muestreo, ensayos a 28 días y curado. & Dosificaciones no presentadas (PAO2, p.~9); instrucciones reiteradas sin registro (pp.~120, 121, 125); ninguna probeta de la losa; columnas y vigas ensayadas a 55--76 días (p.~78). & 0 & ACI~318-25 (pp.~577, 601--603). \\
8.5.2 & Identificar las salidas y su estado de inspección, y asegurar la trazabilidad de materiales y probetas. & Coladas identificadas en los certificados de acero (PAO2, pp.~82--86) sin vínculo con los elementos; probetas rotuladas solo «COLUMNA» o «VIGA» (p.~78). & 0 & --- \\
8.5.3 & Cuidar la propiedad del cliente o de proveedores externos. & --- & NA & La UATF no suministró materiales ni equipos para la losa. \\
8.5.4 & Preservar los materiales y el producto (acopio, protección, curado). & Plan de acopio limitado a «transportar el material y ubicarlo cerca de la obra» (PAO2, p.~10); material y restos dispersos en las fotografías (pp.~75--76); sin registro de curado. & 0 & Cl.~26.9 (PAO2, p.~102). \\
8.5.6 & Revisar y controlar los cambios no planificados en la ejecución. & Modificación estructural y cambio de complemento sin registro (PAO2, pp.~11, 81); vaciado «desde dos frentes» sin plan de juntas (p.~125). & 0 & --- \\
8.6 & Liberar cada etapa solo después de verificar los requisitos, con registro de quién autoriza. & Ocho autorizaciones de vaciado firmadas por la Supervisión (PAO2, pp.~114--125); la losa se certificó para pago el 08/12/2025 (p.~16) antes de que pudiera existir un ensayo a 28 días (26/12/2025). & 0,5 & --- \\
8.7 & Identificar y controlar las salidas no conformes y registrar su tratamiento. & Ningún rechazo, no conformidad ni retrabajo registrado en los 24 asientos (PAO2, pp.~114--125) ni en el informe de Supervisión (pp.~6--11). & 0 & --- \\
%
\multicolumn{5}{@{}l}{\textbf{9. Evaluación del desempeño}}\\
9.1.1 a & Determinar qué se mide (calidad, desperdicio, costos), con qué métodos y cuándo se analizan los resultados. & Solo se mide el avance físico-financiero por planillas (PAO2, pp.~10, 16). & 0,5 & Sin indicadores de calidad ni de desperdicio. \\
9.1.1 b & Medir y conciliar las cantidades ejecutadas (medición conjunta, neta y con respaldo) antes de cada planilla. & Cómputos firmados por contratista, Supervisión y Fiscal (PAO2, pp.~18--32), pero el acero de la losa sin planilla de armado (p.~20), correcciones posteriores (ítem~4: $-$5,66~\mtres{}; acero: $-$1.862,81~kg en la planilla de cierre firmada y $-$1.915,16~kg en su versión digital) y tres cifras distintas del acero acumulado (p.~16; PAO3; Exp., pp.~2, 12). & 0,5 & Cl.~27.ª (PAO2, p.~99). \\
9.1.2 & Hacer seguimiento de la percepción del cliente sobre el cumplimiento de sus requisitos. & El informe de Supervisión califica los trabajos de «buena calidad» y evalúa al Director de Obra (PAO2, pp.~9, 11), sin evidencia objetiva que respalde la calificación. & 0,5 & --- \\
9.1.3 & Analizar y evaluar los datos (conformidad, desempeño de proveedores, eficacia de las acciones). & Sin evidencia; el desfase del periodo se calculó mal: 3,03~\% en lugar de 1,07~\% (PAO2, p.~10). & 0 & --- \\
9.2 & Realizar auditorías internas planificadas. & No se encontraron. & 0 & --- \\
9.3 & Revisar el sistema por la dirección a intervalos planificados. & No se encontró. & 0 & --- \\
%
\multicolumn{5}{@{}l}{\textbf{10. Mejora}}\\
10.1 & Determinar y aplicar oportunidades de mejora. & La Supervisión recomendó presentar las garantías de los materiales, romper probetas y presentar dosificaciones (PAO2, p.~9); no consta que se atendieran en la losa (p.~78). & 0 & --- \\
10.2.1 a & Reaccionar ante las no conformidades: controlarlas, corregirlas y asumir sus consecuencias. & Se corrigieron cantidades después de pagadas: ítem~4, $-$5,66~\mtres{} (PAO2, pp.~16, 18); acero, $-$1.862,81~kg en la planilla de cierre (Exp., p.~2). & 0,5 & Correcciones de medición, sin análisis. \\
10.2.1 b & Analizar las causas, implementar acciones correctivas y verificar su eficacia. & Sin evidencia de análisis de causas ni de acciones correctivas. & 0 & --- \\
10.2.2 & Conservar información documentada de las no conformidades y de las acciones tomadas. & No hay registro de no conformidades. & 0 & --- \\
10.3 & Mejorar continuamente la conveniencia, adecuación y eficacia del sistema. & No se encontró. & 0 & --- \\
\end{longtable}
\vspace{-6pt}
\nota{Elaboración propia. Fuentes: PAO2 = \textcite{uatf2025pao2} (página del archivo PDF; incluye el contrato, el Libro de Órdenes y los certificados); PAO3 = \textcite{uatf_pao3} (archivo de planilla y cómputos de cierre); Exp. = \textcite{uatf_item13} (página del archivo PDF); ET = \textcite{uatf_et}; Plano~02 = sección y detalles de la losa (Anexo de planos). Calificación: 1 = cumple; 0,5 = cumple parcialmente; 0 = no cumple o sin evidencia; NA = no aplica. ACI~318-25 = \textcite{aci318_2025}.}
\endgroup

\textbf{Resumen por cláusula.} La Tabla~\ref{tab:lista-iso-resumen} resume la aplicación. Ningún requisito alcanzó la calificación de cumplimiento pleno: 19 se cumplieron parcialmente y 29 no se cumplieron, de un total de 48 requisitos aplicables.

\begin{table}[H]
\caption{Resumen de la lista de verificación C.1 por cláusula}
\label{tab:lista-iso-resumen}
\setlength{\tabcolsep}{4pt}
\begin{tabularx}{\linewidth}{@{}LQ{1.35cm}Q{1.0cm}Q{1.25cm}Q{1.25cm}Q{1.25cm}Q{1.3cm}Q{1.65cm}@{}}
\toprule
\textbf{Cláusula} & \textbf{Requisitos aplicables} & \textbf{NA} & \textbf{Cumple (1)} & \textbf{Parcial (0,5)} & \textbf{No cumple (0)} & \textbf{Puntaje} & \textbf{Cumplimiento} \\
\midrule
4. Contexto de la organización & 4 & 0 & 0 & 2 & 2 & 1,0 & 25,0~\% \\
5. Liderazgo & 5 & 0 & 0 & 2 & 3 & 1,0 & 20,0~\% \\
6. Planificación & 4 & 0 & 0 & 1 & 3 & 0,5 & 12,5~\% \\
7. Apoyo & 10 & 0 & 0 & 5 & 5 & 2,5 & 25,0~\% \\
8. Operación & 14 & 2 & 0 & 5 & 9 & 2,5 & 17,9~\% \\
9. Evaluación del desempeño & 6 & 0 & 0 & 3 & 3 & 1,5 & 25,0~\% \\
10. Mejora & 5 & 0 & 0 & 1 & 4 & 0,5 & 10,0~\% \\
\midrule
\textbf{Global} & \textbf{48} & \textbf{2} & \textbf{0} & \textbf{19} & \textbf{29} & \textbf{9,5} & \textbf{19,8~\%} \\
\bottomrule
\end{tabularx}
\nota{Elaboración propia. Cumplimiento = puntaje / requisitos aplicables $\times$ 100. El cumplimiento global se calcula con el total de requisitos (9,5/48); el promedio simple de las siete cláusulas es 19,3~\%. Según la escala de \textcite[p.~42]{mamanicalsina2025}, todas las cláusulas y el resultado global se ubican en el rango «deficiente o no implementado» (0--25~\%).}
\end{table}
.
%  Las calificaciones se reproducen en el script
%  03_Datos_Obra_Edificio_Posgrado/04_Extraccion_de_Datos/scripts/
%  cap2_diagnostico_figuras.py (Figura de brechas del Capítulo 2).
% =====================================================================
\section*{C.1. Lista de verificación de requisitos de ISO~9001:2026 aplicada a la ejecución de la primera losa}
\addcontentsline{toc}{section}{C.1. Lista de verificación de requisitos de ISO~9001:2026 (primera losa)}
\label{anx:lista-iso}

\textbf{Propósito y alcance.} La lista sirve para diagnosticar las brechas entre la gestión real del proceso constructivo de la losa alivianada y los requisitos de las cláusulas 4 a 10 de la norma ISO~9001:2026 \parencite{iso9001_2026}. Se aplicó a la ejecución de la primera losa (41 paños; 1.109,84~\mdos{}), entendida como el sistema formado por la empresa contratista, su Dirección de Obra y los controles de la Supervisión que constan en los documentos. Los requisitos se redactaron como paráfrasis del autor, en términos del proceso de la losa; no reproducen el texto de la norma. La numeración corresponde a la estructura de cláusulas que la edición de 2026 conserva de la de 2015. \pendiente{contrastar la numeración de las subcláusulas con el texto de la NB/ISO~9001:2026 de IBNORCA \parencite{ibnorca9001_2026}.}

\textbf{Método.} La aplicación fue documental y la realizó el investigador. Se revisaron, página por página, la Planilla de avance de obra N.\textsuperscript{o}~2 con sus anexos (informe de Supervisión, cómputos, certificados, contrato, Libro de Órdenes y cronogramas), la planilla y los cómputos de cierre por resolución del contrato, el expediente del ítem~13, las especificaciones técnicas y los planos 01 a 03. Cada requisito aplicable se calificó así: \textbf{1} = cumple (el requisito está establecido, se aplicó y hay evidencia documental); \textbf{0,5} = cumple parcialmente (hay evidencia de una aplicación parcial, informal o sin registros completos); \textbf{0} = no cumple (no se encontró evidencia en los documentos revisados, o la evidencia muestra lo contrario); \textbf{NA} = no aplica al proceso de la losa (se excluye del denominador). El porcentaje de cada cláusula es la suma de las calificaciones dividida entre el número de requisitos aplicables. La interpretación usa la escala de \textcite[p.~42]{mamanicalsina2025}: 0--25~\% deficiente o no implementado; 26--50~\% básico o en desarrollo; 51--75~\% aceptable; 76--100~\% excelente.

\textbf{Limitación.} Una calificación de 0 por falta de evidencia documental no prueba que la práctica no existiera; prueba que no quedó registrada. Las entrevistas de la guía C.7 a la Dirección de Obra, a la Supervisión y a la Fiscalización permitirán verificar o corregir estas calificaciones. Para las auditorías internas del SGC propuesto, la columna de requisitos se usa sin cambios y las columnas de evidencia y calificación se llenan en cada aplicación.

\begingroup\footnotesize\setstretch{1}\setlength{\tabcolsep}{3pt}
\begin{longtable}{@{}P{1.15cm}P{3.85cm}P{5.05cm}Q{0.85cm}P{4.15cm}@{}}
\caption{Lista de verificación C.1 aplicada a la ejecución de la primera losa}\label{tab:lista-iso}\\
\toprule
\textbf{Cláusula} & \textbf{Requisito (paráfrasis aplicada a la losa)} & \textbf{Evidencia documental encontrada} & \textbf{Calif.} & \textbf{Observación} \\
\midrule
\endfirsthead
\multicolumn{5}{@{}l}{\textit{(continuación)}}\\
\toprule
\textbf{Cláusula} & \textbf{Requisito (paráfrasis aplicada a la losa)} & \textbf{Evidencia documental encontrada} & \textbf{Calif.} & \textbf{Observación} \\
\midrule
\endhead
\midrule
\multicolumn{5}{@{}r}{\textit{(continúa)}}\\
\endfoot
\bottomrule
\endlastfoot
%
\multicolumn{5}{@{}l}{\textbf{4. Contexto de la organización}}\\
4.1 & Determinar las cuestiones externas e internas que afectan al proceso de la losa (abastecimiento, precios, clima y pertinencia del cambio climático) y darles seguimiento. & No hay análisis de contexto. Las lluvias (siete días) y la demora de materiales por escasez de combustible solo se anotaron después de ocurridas, para pedir ampliación de plazo (PAO2, pp.~114--124). & 0 & La escasez de diésel ya era pública en junio de 2025 (FANCESA; Las Lomas) y no generó previsiones. \\
4.2 & Identificar las partes interesadas pertinentes (UATF, Supervisión, Fiscalización, proveedores, laboratorio, usuarios) y sus requisitos. & El contrato define partes y obligaciones (PAO2, pp.~91--112); la ET y los planos fijan los requisitos del producto. No hay análisis de las demás partes interesadas. & 0,5 & Documentados solo los requisitos de la entidad contratante. \\
4.3 & Definir y documentar el alcance del sistema, con los requisitos que no se aplican y su justificación. & No existe un sistema declarado ni un documento de alcance; los 15 documentos de la planilla no incluyen un plan de la calidad (PAO2, pp.~3--4, 12). & 0 & --- \\
4.4 & Determinar los procesos de la losa (abastecimiento, montaje, armado, hormigonado, curado y medición) con entradas, salidas, secuencia, criterios, responsables y recursos. & La secuencia general figura en la ET (pp.~1--4) y en el cronograma A-8 (PAO2, p.~137); no hay mapa de procesos, procedimientos ni responsables por actividad. & 0,5 & El cronograma no refleja la secuencia real: la losa se armó en noviembre de 2025. \\
%
\multicolumn{5}{@{}l}{\textbf{5. Liderazgo}}\\
5.1.1 a & Demostrar el compromiso de la dirección de la empresa con la eficacia del sistema (rendición de cuentas, recursos, integración en los procesos). & La dirección de la empresa solo aparece en las cartas de entrega de planillas y en las firmas de cobro (PAO2, pp.~12, 14, 88). & 0 & --- \\
5.1.1 b & Promover una cultura de la calidad y un comportamiento ético en la ejecución, la medición y el registro de la obra. & No hay evidencia de acciones; las diferencias de cantidades del acero no tienen explicación documentada (PAO3, hoja de cómputos, filas 375--541; Exp., p.~12). & 0 & Énfasis de la edición 2026. Es ausencia de evidencia, no un juicio sobre las personas. \\
5.1.2 & Mantener el enfoque al cliente: cumplir los requisitos de la UATF y los legales, y atender los riesgos que afectan la conformidad de la losa. & Obra ejecutada con autorización de la Supervisión (PAO2, pp.~114--125), pero con retraso acumulado de 3,03~\% (p.~10) y con garantías, ensayos y dosificaciones pendientes al cierre del periodo (p.~9). & 0,5 & --- \\
5.2 & Establecer, comunicar y mantener disponible una política de la calidad apropiada al proyecto. & No se encontró. & 0 & --- \\
5.3 & Asignar y comunicar roles, responsabilidades y autoridades, incluido quién informa sobre el desempeño del sistema. & Director de Obra designado y evaluado (PAO2, p.~9); Supervisora y Fiscal designados (pp.~7, 127). No hay organigrama ni responsable de calidad; falta la constancia del Director de Obra en las 12 hojas del Libro de Órdenes (pp.~114--125). & 0,5 & Cl.~25.ª del contrato (PAO2, p.~100). \\
%
\multicolumn{5}{@{}l}{\textbf{6. Planificación}}\\
6.1 a & Determinar los riesgos y las oportunidades del proceso (suministro, precios, clima, cambios de diseño). & No hay matriz de riesgos. Se materializaron riesgos no previstos: viguetas inexistentes en el mercado en los tamaños requeridos (PAO2, p.~11), entrega del 30~\% del acero pedido (carta de Las Lomas, Exp., p.~27) y demoras por combustible (PAO2, p.~123). & 0 & --- \\
6.1 b & Planificar acciones proporcionales a los riesgos, integrarlas en los procesos y evaluar su eficacia. & Sin evidencia; los eventos solo se tramitaron como causales de ampliación de plazo (PAO2, pp.~114--124), aunque la cl.~20.ª no los considera fuerza mayor (p.~96). & 0 & --- \\
6.2 & Establecer objetivos de la calidad medibles (desperdicio, conformidad del hormigón, plazos de liberación) y planificar cómo lograrlos. & El único objetivo medido es el avance físico-financiero (PAO2, p.~10). & 0 & --- \\
6.3 & Planificar y controlar los cambios que afectan al sistema o al producto. & Los cambios contractuales se tramitaron formalmente: contrato modificatorio N.\textsuperscript{o}~1 y dos ampliaciones de plazo (PAO2, pp.~4, 7). Los cambios técnicos de la losa (viguetas, complemento) no tienen registro (p.~11; p.~81; ET, p.~3). & 0,5 & --- \\
%
\multicolumn{5}{@{}l}{\textbf{7. Apoyo}}\\
7.1.2 y 7.1.3 & Proveer las personas y la infraestructura necesarias (frentes de trabajo, apuntalamiento, mezclado, transporte y vibrado del hormigón). & Dos frentes de obra civil y dos de vaciado de $\approx$~6 trabajadores (PAO2, p.~9); apuntalamiento de madera, carretillas y vibrador en las fotografías (pp.~74--76). La Supervisión pidió más cuadrillas para recuperar el retraso (p.~11). & 0,5 & Recursos presentes, pero insuficientes para el cronograma. \\
7.1.4 & Mantener un ambiente adecuado para la operación (protección del hormigón ante lluvia y bajas temperaturas). & No hay previsiones de vaciado ni de curado en tiempo frío o con lluvia; solo el registro de las lluvias (PAO2, pp.~114--124). & 0 & ACI~318-25 exige un plan de hormigonado en tiempo frío (p.~583). \\
7.1.5 & Disponer de recursos de seguimiento y medición adecuados y de un muestreo confiable (laboratorio, moldes, cono de Abrams, balanzas). & Ensayos en el laboratorio de la UATF (PAO2, p.~78), con probetas «traídas y elaboradas por la empresa solicitante»; no constan equipos de obra ni su verificación. & 0,5 & --- \\
7.1.6 & Determinar y mantener los conocimientos necesarios (instrucciones de montaje del fabricante, lecciones aprendidas). & El certificado de PRETENSA remite a los pasos de montaje del fabricante (PAO2, p.~80); no consta que estuvieran en obra. & 0 & --- \\
7.2 & Determinar la competencia necesaria (armadores, encofradores, personal de vaciado y muestreo), asegurarla y conservar evidencia. & Solo se evalúa al Director de Obra («Buena», PAO2, p.~9); la Supervisión trabajó con apoyo de pasantes (p.~10); no hay registros de competencia ni de capacitación. & 0,5 & --- \\
7.3 & Asegurar que el personal conozca la importancia de su trabajo para la calidad y las consecuencias de incumplir. & Sin evidencia; las fotografías muestran la parrilla pisada durante el vaciado y sin separadores visibles (PAO2, p.~76). & 0 & Observación fotográfica a confirmar. \\
7.4 & Definir la comunicación interna y externa (qué, cuándo, a quién, cómo y quién comunica). & El Libro de Órdenes funcionó como canal formal: 24 asientos con respuesta el mismo día (PAO2, pp.~114--125); hay cartas formales (pp.~2, 5, 12, 88). Encabezados vacíos y sin constancia del Director de Obra. & 0,5 & Sin evidencia de comunicación interna con el personal. \\
7.5.2 & Crear y actualizar la información documentada con identificación, fecha, revisión y aprobación. & Planillas, cómputos y croquis firmados y sellados (PAO2, pp.~14, 18--72), pero con errores: rótulo ZT10 repetido (p.~18), columna tipo~5 con 6 y 7 unidades (p.~38), hoja duplicada (pp.~51--52), anticipo de 10~\% en un contrato sin anticipo (p.~14), desfase del periodo mal calculado (p.~10). & 0,5 & --- \\
7.5.3 a & Controlar la información documentada: versión vigente, coherencia entre documentos y normas actualizadas. & Tres geometrías del complemento: 100~$\times$~50~$\times$~20~cm (ET, p.~3), altura de 0,15~m (Plano~02) y 20~$\times$~100~$\times$~44~cm (PAO2, p.~81); la ET cita la CBH-87 y NB antiguas (pp.~1--2); dos versiones de la planilla de cierre con distinto descuento de acero (Exp., p.~2; PAO3); el archivo de cierre conserva una hoja de otro contrato (PAO3). & 0 & --- \\
7.5.3 b & Conservar los registros que evidencian la conformidad (recepción, inspección, ensayos, consumos, vaciado y curado). & No se encontraron registros de recepción, inspección del armado, consumos, vaciado ni curado de la losa (PAO2, pp.~1--139; Exp., pp.~1--29). & 0 & Brecha central para la línea base. \\
%
\multicolumn{5}{@{}l}{\textbf{8. Operación}}\\
8.1 a & Planificar el proceso de la losa: criterios de aceptación, controles y registros necesarios (plan de la calidad, PPI). & La ET fija criterios: f'c~=~250~kg/cm\textsuperscript{2}, cemento $\geq$~400~kg/\mtres{}, vaciado por franjas y monolítico con las vigas (pp.~1--4). No hay plan de la calidad ni PPI. & 0,5 & --- \\
8.1 b & Determinar los materiales necesarios por paño (cantidades teóricas para pedido y control). & El cómputo del ítem~13 solo lista las áreas de los 41 paños (PAO2, pp.~28--29); no hay cantidades de viguetas, EPS, hormigón ni acero por paño. & 0 & --- \\
8.2.2 y 8.2.3 & Determinar y revisar los requisitos del producto (contrato, ET, planos) antes de ejecutar, y resolver las diferencias. & Requisitos disponibles (PAO2, pp.~91--112; ET; planos 01 a 03), pero las discrepancias del complemento, de su altura y de las normas citadas no se resolvieron, pese a la consulta previa que exige la cl.~33.1 (p.~107). & 0,5 & --- \\
8.2.4 & Actualizar la información y comunicarla cuando cambian los requisitos. & «Modificaciones estructurales» por falta de viguetas de los tamaños requeridos, sin documento que las describa (PAO2, p.~11); EPS distinto de la ET sin orden de cambio (p.~81). & 0 & --- \\
8.3 & Diseño y desarrollo de los productos. & --- & NA & El diseño estructural corresponde al proyectista de la UATF. \\
8.4.1 & Aplicar criterios para evaluar, seleccionar y hacer seguimiento a los proveedores externos. & Sin criterios documentados ni reevaluación tras los incumplimientos (acero entregado en 30~\%, Exp., p.~27; viguetas no disponibles, PAO2, p.~11). La certificación NB~997 e ISO~9001 de PRETENSA (p.~80) es un atributo del proveedor, no un criterio registrado. & 0 & --- \\
8.4.2 & Verificar lo suministrado antes de incorporarlo a la obra (certificados, inspección de recepción, ensayos). & Certificados de PRETENSA y UNICON del 27/11/2025, sin firma y posteriores al inicio del armado (14/11/2025) (PAO2, pp.~80--81, 122); certificados de acero de $\phi$20 a $\phi$8 (pp.~82--86), sin $\phi$6 ni $\phi$10; sin certificados de cemento ni áridos; sin inspección de recepción. & 0,5 & Cl.~26.6 (PAO2, p.~102); ACI~318-25 (p.~587). \\
8.4.3 & Comunicar a los proveedores los requisitos de lo que se compra (tipo, dimensiones, cantidad, aprobación). & No hay órdenes de compra; el EPS entregado (20~$\times$~100~$\times$~44~cm; 2.350 piezas) no coincide con la dimensión de la ET (PAO2, p.~81; ET, p.~3). & 0 & --- \\
8.5.1 a & Ejecutar en condiciones controladas: información disponible, seguimiento en las etapas apropiadas y personal competente. & Verificación del armado antes de cada vaciado, registrada en el Libro de Órdenes (PAO2, pp.~114--125), sin lista de verificación ni criterios; no consta la fecha del hormigonado de la losa (p.~125). & 0,5 & --- \\
8.5.1 b & Controlar el hormigón y validar el proceso de hormigonado: dosificación aprobada, asentamiento, muestreo, ensayos a 28 días y curado. & Dosificaciones no presentadas (PAO2, p.~9); instrucciones reiteradas sin registro (pp.~120, 121, 125); ninguna probeta de la losa; columnas y vigas ensayadas a 55--76 días (p.~78). & 0 & ACI~318-25 (pp.~577, 601--603). \\
8.5.2 & Identificar las salidas y su estado de inspección, y asegurar la trazabilidad de materiales y probetas. & Coladas identificadas en los certificados de acero (PAO2, pp.~82--86) sin vínculo con los elementos; probetas rotuladas solo «COLUMNA» o «VIGA» (p.~78). & 0 & --- \\
8.5.3 & Cuidar la propiedad del cliente o de proveedores externos. & --- & NA & La UATF no suministró materiales ni equipos para la losa. \\
8.5.4 & Preservar los materiales y el producto (acopio, protección, curado). & Plan de acopio limitado a «transportar el material y ubicarlo cerca de la obra» (PAO2, p.~10); material y restos dispersos en las fotografías (pp.~75--76); sin registro de curado. & 0 & Cl.~26.9 (PAO2, p.~102). \\
8.5.6 & Revisar y controlar los cambios no planificados en la ejecución. & Modificación estructural y cambio de complemento sin registro (PAO2, pp.~11, 81); vaciado «desde dos frentes» sin plan de juntas (p.~125). & 0 & --- \\
8.6 & Liberar cada etapa solo después de verificar los requisitos, con registro de quién autoriza. & Ocho autorizaciones de vaciado firmadas por la Supervisión (PAO2, pp.~114--125); la losa se certificó para pago el 08/12/2025 (p.~16) antes de que pudiera existir un ensayo a 28 días (26/12/2025). & 0,5 & --- \\
8.7 & Identificar y controlar las salidas no conformes y registrar su tratamiento. & Ningún rechazo, no conformidad ni retrabajo registrado en los 24 asientos (PAO2, pp.~114--125) ni en el informe de Supervisión (pp.~6--11). & 0 & --- \\
%
\multicolumn{5}{@{}l}{\textbf{9. Evaluación del desempeño}}\\
9.1.1 a & Determinar qué se mide (calidad, desperdicio, costos), con qué métodos y cuándo se analizan los resultados. & Solo se mide el avance físico-financiero por planillas (PAO2, pp.~10, 16). & 0,5 & Sin indicadores de calidad ni de desperdicio. \\
9.1.1 b & Medir y conciliar las cantidades ejecutadas (medición conjunta, neta y con respaldo) antes de cada planilla. & Cómputos firmados por contratista, Supervisión y Fiscal (PAO2, pp.~18--32), pero el acero de la losa sin planilla de armado (p.~20), correcciones posteriores (ítem~4: $-$5,66~\mtres{}; acero: $-$1.862,81~kg en la planilla de cierre firmada y $-$1.915,16~kg en su versión digital) y tres cifras distintas del acero acumulado (p.~16; PAO3; Exp., pp.~2, 12). & 0,5 & Cl.~27.ª (PAO2, p.~99). \\
9.1.2 & Hacer seguimiento de la percepción del cliente sobre el cumplimiento de sus requisitos. & El informe de Supervisión califica los trabajos de «buena calidad» y evalúa al Director de Obra (PAO2, pp.~9, 11), sin evidencia objetiva que respalde la calificación. & 0,5 & --- \\
9.1.3 & Analizar y evaluar los datos (conformidad, desempeño de proveedores, eficacia de las acciones). & Sin evidencia; el desfase del periodo se calculó mal: 3,03~\% en lugar de 1,07~\% (PAO2, p.~10). & 0 & --- \\
9.2 & Realizar auditorías internas planificadas. & No se encontraron. & 0 & --- \\
9.3 & Revisar el sistema por la dirección a intervalos planificados. & No se encontró. & 0 & --- \\
%
\multicolumn{5}{@{}l}{\textbf{10. Mejora}}\\
10.1 & Determinar y aplicar oportunidades de mejora. & La Supervisión recomendó presentar las garantías de los materiales, romper probetas y presentar dosificaciones (PAO2, p.~9); no consta que se atendieran en la losa (p.~78). & 0 & --- \\
10.2.1 a & Reaccionar ante las no conformidades: controlarlas, corregirlas y asumir sus consecuencias. & Se corrigieron cantidades después de pagadas: ítem~4, $-$5,66~\mtres{} (PAO2, pp.~16, 18); acero, $-$1.862,81~kg en la planilla de cierre (Exp., p.~2). & 0,5 & Correcciones de medición, sin análisis. \\
10.2.1 b & Analizar las causas, implementar acciones correctivas y verificar su eficacia. & Sin evidencia de análisis de causas ni de acciones correctivas. & 0 & --- \\
10.2.2 & Conservar información documentada de las no conformidades y de las acciones tomadas. & No hay registro de no conformidades. & 0 & --- \\
10.3 & Mejorar continuamente la conveniencia, adecuación y eficacia del sistema. & No se encontró. & 0 & --- \\
\end{longtable}
\vspace{-6pt}
\nota{Elaboración propia. Fuentes: PAO2 = \textcite{uatf2025pao2} (página del archivo PDF; incluye el contrato, el Libro de Órdenes y los certificados); PAO3 = \textcite{uatf_pao3} (archivo de planilla y cómputos de cierre); Exp. = \textcite{uatf_item13} (página del archivo PDF); ET = \textcite{uatf_et}; Plano~02 = sección y detalles de la losa (Anexo de planos). Calificación: 1 = cumple; 0,5 = cumple parcialmente; 0 = no cumple o sin evidencia; NA = no aplica. ACI~318-25 = \textcite{aci318_2025}.}
\endgroup

\textbf{Resumen por cláusula.} La Tabla~\ref{tab:lista-iso-resumen} resume la aplicación. Ningún requisito alcanzó la calificación de cumplimiento pleno: 19 se cumplieron parcialmente y 29 no se cumplieron, de un total de 48 requisitos aplicables.

\begin{table}[H]
\caption{Resumen de la lista de verificación C.1 por cláusula}
\label{tab:lista-iso-resumen}
\setlength{\tabcolsep}{4pt}
\begin{tabularx}{\linewidth}{@{}LQ{1.35cm}Q{1.0cm}Q{1.25cm}Q{1.25cm}Q{1.25cm}Q{1.3cm}Q{1.65cm}@{}}
\toprule
\textbf{Cláusula} & \textbf{Requisitos aplicables} & \textbf{NA} & \textbf{Cumple (1)} & \textbf{Parcial (0,5)} & \textbf{No cumple (0)} & \textbf{Puntaje} & \textbf{Cumplimiento} \\
\midrule
4. Contexto de la organización & 4 & 0 & 0 & 2 & 2 & 1,0 & 25,0~\% \\
5. Liderazgo & 5 & 0 & 0 & 2 & 3 & 1,0 & 20,0~\% \\
6. Planificación & 4 & 0 & 0 & 1 & 3 & 0,5 & 12,5~\% \\
7. Apoyo & 10 & 0 & 0 & 5 & 5 & 2,5 & 25,0~\% \\
8. Operación & 14 & 2 & 0 & 5 & 9 & 2,5 & 17,9~\% \\
9. Evaluación del desempeño & 6 & 0 & 0 & 3 & 3 & 1,5 & 25,0~\% \\
10. Mejora & 5 & 0 & 0 & 1 & 4 & 0,5 & 10,0~\% \\
\midrule
\textbf{Global} & \textbf{48} & \textbf{2} & \textbf{0} & \textbf{19} & \textbf{29} & \textbf{9,5} & \textbf{19,8~\%} \\
\bottomrule
\end{tabularx}
\nota{Elaboración propia. Cumplimiento = puntaje / requisitos aplicables $\times$ 100. El cumplimiento global se calcula con el total de requisitos (9,5/48); el promedio simple de las siete cláusulas es 19,3~\%. Según la escala de \textcite[p.~42]{mamanicalsina2025}, todas las cláusulas y el resultado global se ubican en el rango «deficiente o no implementado» (0--25~\%).}
\end{table}
.
%  Las calificaciones se reproducen en el script
%  03_Datos_Obra_Edificio_Posgrado/04_Extraccion_de_Datos/scripts/
%  cap2_diagnostico_figuras.py (Figura de brechas del Capítulo 2).
% =====================================================================
\section*{C.1. Lista de verificación de requisitos de ISO~9001:2026 aplicada a la ejecución de la primera losa}
\addcontentsline{toc}{section}{C.1. Lista de verificación de requisitos de ISO~9001:2026 (primera losa)}
\label{anx:lista-iso}

\textbf{Propósito y alcance.} La lista sirve para diagnosticar las brechas entre la gestión real del proceso constructivo de la losa alivianada y los requisitos de las cláusulas 4 a 10 de la norma ISO~9001:2026 \parencite{iso9001_2026}. Se aplicó a la ejecución de la primera losa (41 paños; 1.109,84~\mdos{}), entendida como el sistema formado por la empresa contratista, su Dirección de Obra y los controles de la Supervisión que constan en los documentos. Los requisitos se redactaron como paráfrasis del autor, en términos del proceso de la losa; no reproducen el texto de la norma. La numeración corresponde a la estructura de cláusulas que la edición de 2026 conserva de la de 2015. \pendiente{contrastar la numeración de las subcláusulas con el texto de la NB/ISO~9001:2026 de IBNORCA \parencite{ibnorca9001_2026}.}

\textbf{Método.} La aplicación fue documental y la realizó el investigador. Se revisaron, página por página, la Planilla de avance de obra N.\textsuperscript{o}~2 con sus anexos (informe de Supervisión, cómputos, certificados, contrato, Libro de Órdenes y cronogramas), la planilla y los cómputos de cierre por resolución del contrato, el expediente del ítem~13, las especificaciones técnicas y los planos 01 a 03. Cada requisito aplicable se calificó así: \textbf{1} = cumple (el requisito está establecido, se aplicó y hay evidencia documental); \textbf{0,5} = cumple parcialmente (hay evidencia de una aplicación parcial, informal o sin registros completos); \textbf{0} = no cumple (no se encontró evidencia en los documentos revisados, o la evidencia muestra lo contrario); \textbf{NA} = no aplica al proceso de la losa (se excluye del denominador). El porcentaje de cada cláusula es la suma de las calificaciones dividida entre el número de requisitos aplicables. La interpretación usa la escala de \textcite[p.~42]{mamanicalsina2025}: 0--25~\% deficiente o no implementado; 26--50~\% básico o en desarrollo; 51--75~\% aceptable; 76--100~\% excelente.

\textbf{Limitación.} Una calificación de 0 por falta de evidencia documental no prueba que la práctica no existiera; prueba que no quedó registrada. Las entrevistas de la guía C.7 a la Dirección de Obra, a la Supervisión y a la Fiscalización permitirán verificar o corregir estas calificaciones. Para las auditorías internas del SGC propuesto, la columna de requisitos se usa sin cambios y las columnas de evidencia y calificación se llenan en cada aplicación.

\begingroup\footnotesize\setstretch{1}\setlength{\tabcolsep}{3pt}
\begin{longtable}{@{}P{1.15cm}P{3.85cm}P{5.05cm}Q{0.85cm}P{4.15cm}@{}}
\caption{Lista de verificación C.1 aplicada a la ejecución de la primera losa}\label{tab:lista-iso}\\
\toprule
\textbf{Cláusula} & \textbf{Requisito (paráfrasis aplicada a la losa)} & \textbf{Evidencia documental encontrada} & \textbf{Calif.} & \textbf{Observación} \\
\midrule
\endfirsthead
\multicolumn{5}{@{}l}{\textit{(continuación)}}\\
\toprule
\textbf{Cláusula} & \textbf{Requisito (paráfrasis aplicada a la losa)} & \textbf{Evidencia documental encontrada} & \textbf{Calif.} & \textbf{Observación} \\
\midrule
\endhead
\midrule
\multicolumn{5}{@{}r}{\textit{(continúa)}}\\
\endfoot
\bottomrule
\endlastfoot
%
\multicolumn{5}{@{}l}{\textbf{4. Contexto de la organización}}\\
4.1 & Determinar las cuestiones externas e internas que afectan al proceso de la losa (abastecimiento, precios, clima y pertinencia del cambio climático) y darles seguimiento. & No hay análisis de contexto. Las lluvias (siete días) y la demora de materiales por escasez de combustible solo se anotaron después de ocurridas, para pedir ampliación de plazo (PAO2, pp.~114--124). & 0 & La escasez de diésel ya era pública en junio de 2025 (FANCESA; Las Lomas) y no generó previsiones. \\
4.2 & Identificar las partes interesadas pertinentes (UATF, Supervisión, Fiscalización, proveedores, laboratorio, usuarios) y sus requisitos. & El contrato define partes y obligaciones (PAO2, pp.~91--112); la ET y los planos fijan los requisitos del producto. No hay análisis de las demás partes interesadas. & 0,5 & Documentados solo los requisitos de la entidad contratante. \\
4.3 & Definir y documentar el alcance del sistema, con los requisitos que no se aplican y su justificación. & No existe un sistema declarado ni un documento de alcance; los 15 documentos de la planilla no incluyen un plan de la calidad (PAO2, pp.~3--4, 12). & 0 & --- \\
4.4 & Determinar los procesos de la losa (abastecimiento, montaje, armado, hormigonado, curado y medición) con entradas, salidas, secuencia, criterios, responsables y recursos. & La secuencia general figura en la ET (pp.~1--4) y en el cronograma A-8 (PAO2, p.~137); no hay mapa de procesos, procedimientos ni responsables por actividad. & 0,5 & El cronograma no refleja la secuencia real: la losa se armó en noviembre de 2025. \\
%
\multicolumn{5}{@{}l}{\textbf{5. Liderazgo}}\\
5.1.1 a & Demostrar el compromiso de la dirección de la empresa con la eficacia del sistema (rendición de cuentas, recursos, integración en los procesos). & La dirección de la empresa solo aparece en las cartas de entrega de planillas y en las firmas de cobro (PAO2, pp.~12, 14, 88). & 0 & --- \\
5.1.1 b & Promover una cultura de la calidad y un comportamiento ético en la ejecución, la medición y el registro de la obra. & No hay evidencia de acciones; las diferencias de cantidades del acero no tienen explicación documentada (PAO3, hoja de cómputos, filas 375--541; Exp., p.~12). & 0 & Énfasis de la edición 2026. Es ausencia de evidencia, no un juicio sobre las personas. \\
5.1.2 & Mantener el enfoque al cliente: cumplir los requisitos de la UATF y los legales, y atender los riesgos que afectan la conformidad de la losa. & Obra ejecutada con autorización de la Supervisión (PAO2, pp.~114--125), pero con retraso acumulado de 3,03~\% (p.~10) y con garantías, ensayos y dosificaciones pendientes al cierre del periodo (p.~9). & 0,5 & --- \\
5.2 & Establecer, comunicar y mantener disponible una política de la calidad apropiada al proyecto. & No se encontró. & 0 & --- \\
5.3 & Asignar y comunicar roles, responsabilidades y autoridades, incluido quién informa sobre el desempeño del sistema. & Director de Obra designado y evaluado (PAO2, p.~9); Supervisora y Fiscal designados (pp.~7, 127). No hay organigrama ni responsable de calidad; falta la constancia del Director de Obra en las 12 hojas del Libro de Órdenes (pp.~114--125). & 0,5 & Cl.~25.ª del contrato (PAO2, p.~100). \\
%
\multicolumn{5}{@{}l}{\textbf{6. Planificación}}\\
6.1 a & Determinar los riesgos y las oportunidades del proceso (suministro, precios, clima, cambios de diseño). & No hay matriz de riesgos. Se materializaron riesgos no previstos: viguetas inexistentes en el mercado en los tamaños requeridos (PAO2, p.~11), entrega del 30~\% del acero pedido (carta de Las Lomas, Exp., p.~27) y demoras por combustible (PAO2, p.~123). & 0 & --- \\
6.1 b & Planificar acciones proporcionales a los riesgos, integrarlas en los procesos y evaluar su eficacia. & Sin evidencia; los eventos solo se tramitaron como causales de ampliación de plazo (PAO2, pp.~114--124), aunque la cl.~20.ª no los considera fuerza mayor (p.~96). & 0 & --- \\
6.2 & Establecer objetivos de la calidad medibles (desperdicio, conformidad del hormigón, plazos de liberación) y planificar cómo lograrlos. & El único objetivo medido es el avance físico-financiero (PAO2, p.~10). & 0 & --- \\
6.3 & Planificar y controlar los cambios que afectan al sistema o al producto. & Los cambios contractuales se tramitaron formalmente: contrato modificatorio N.\textsuperscript{o}~1 y dos ampliaciones de plazo (PAO2, pp.~4, 7). Los cambios técnicos de la losa (viguetas, complemento) no tienen registro (p.~11; p.~81; ET, p.~3). & 0,5 & --- \\
%
\multicolumn{5}{@{}l}{\textbf{7. Apoyo}}\\
7.1.2 y 7.1.3 & Proveer las personas y la infraestructura necesarias (frentes de trabajo, apuntalamiento, mezclado, transporte y vibrado del hormigón). & Dos frentes de obra civil y dos de vaciado de $\approx$~6 trabajadores (PAO2, p.~9); apuntalamiento de madera, carretillas y vibrador en las fotografías (pp.~74--76). La Supervisión pidió más cuadrillas para recuperar el retraso (p.~11). & 0,5 & Recursos presentes, pero insuficientes para el cronograma. \\
7.1.4 & Mantener un ambiente adecuado para la operación (protección del hormigón ante lluvia y bajas temperaturas). & No hay previsiones de vaciado ni de curado en tiempo frío o con lluvia; solo el registro de las lluvias (PAO2, pp.~114--124). & 0 & ACI~318-25 exige un plan de hormigonado en tiempo frío (p.~583). \\
7.1.5 & Disponer de recursos de seguimiento y medición adecuados y de un muestreo confiable (laboratorio, moldes, cono de Abrams, balanzas). & Ensayos en el laboratorio de la UATF (PAO2, p.~78), con probetas «traídas y elaboradas por la empresa solicitante»; no constan equipos de obra ni su verificación. & 0,5 & --- \\
7.1.6 & Determinar y mantener los conocimientos necesarios (instrucciones de montaje del fabricante, lecciones aprendidas). & El certificado de PRETENSA remite a los pasos de montaje del fabricante (PAO2, p.~80); no consta que estuvieran en obra. & 0 & --- \\
7.2 & Determinar la competencia necesaria (armadores, encofradores, personal de vaciado y muestreo), asegurarla y conservar evidencia. & Solo se evalúa al Director de Obra («Buena», PAO2, p.~9); la Supervisión trabajó con apoyo de pasantes (p.~10); no hay registros de competencia ni de capacitación. & 0,5 & --- \\
7.3 & Asegurar que el personal conozca la importancia de su trabajo para la calidad y las consecuencias de incumplir. & Sin evidencia; las fotografías muestran la parrilla pisada durante el vaciado y sin separadores visibles (PAO2, p.~76). & 0 & Observación fotográfica a confirmar. \\
7.4 & Definir la comunicación interna y externa (qué, cuándo, a quién, cómo y quién comunica). & El Libro de Órdenes funcionó como canal formal: 24 asientos con respuesta el mismo día (PAO2, pp.~114--125); hay cartas formales (pp.~2, 5, 12, 88). Encabezados vacíos y sin constancia del Director de Obra. & 0,5 & Sin evidencia de comunicación interna con el personal. \\
7.5.2 & Crear y actualizar la información documentada con identificación, fecha, revisión y aprobación. & Planillas, cómputos y croquis firmados y sellados (PAO2, pp.~14, 18--72), pero con errores: rótulo ZT10 repetido (p.~18), columna tipo~5 con 6 y 7 unidades (p.~38), hoja duplicada (pp.~51--52), anticipo de 10~\% en un contrato sin anticipo (p.~14), desfase del periodo mal calculado (p.~10). & 0,5 & --- \\
7.5.3 a & Controlar la información documentada: versión vigente, coherencia entre documentos y normas actualizadas. & Tres geometrías del complemento: 100~$\times$~50~$\times$~20~cm (ET, p.~3), altura de 0,15~m (Plano~02) y 20~$\times$~100~$\times$~44~cm (PAO2, p.~81); la ET cita la CBH-87 y NB antiguas (pp.~1--2); dos versiones de la planilla de cierre con distinto descuento de acero (Exp., p.~2; PAO3); el archivo de cierre conserva una hoja de otro contrato (PAO3). & 0 & --- \\
7.5.3 b & Conservar los registros que evidencian la conformidad (recepción, inspección, ensayos, consumos, vaciado y curado). & No se encontraron registros de recepción, inspección del armado, consumos, vaciado ni curado de la losa (PAO2, pp.~1--139; Exp., pp.~1--29). & 0 & Brecha central para la línea base. \\
%
\multicolumn{5}{@{}l}{\textbf{8. Operación}}\\
8.1 a & Planificar el proceso de la losa: criterios de aceptación, controles y registros necesarios (plan de la calidad, PPI). & La ET fija criterios: f'c~=~250~kg/cm\textsuperscript{2}, cemento $\geq$~400~kg/\mtres{}, vaciado por franjas y monolítico con las vigas (pp.~1--4). No hay plan de la calidad ni PPI. & 0,5 & --- \\
8.1 b & Determinar los materiales necesarios por paño (cantidades teóricas para pedido y control). & El cómputo del ítem~13 solo lista las áreas de los 41 paños (PAO2, pp.~28--29); no hay cantidades de viguetas, EPS, hormigón ni acero por paño. & 0 & --- \\
8.2.2 y 8.2.3 & Determinar y revisar los requisitos del producto (contrato, ET, planos) antes de ejecutar, y resolver las diferencias. & Requisitos disponibles (PAO2, pp.~91--112; ET; planos 01 a 03), pero las discrepancias del complemento, de su altura y de las normas citadas no se resolvieron, pese a la consulta previa que exige la cl.~33.1 (p.~107). & 0,5 & --- \\
8.2.4 & Actualizar la información y comunicarla cuando cambian los requisitos. & «Modificaciones estructurales» por falta de viguetas de los tamaños requeridos, sin documento que las describa (PAO2, p.~11); EPS distinto de la ET sin orden de cambio (p.~81). & 0 & --- \\
8.3 & Diseño y desarrollo de los productos. & --- & NA & El diseño estructural corresponde al proyectista de la UATF. \\
8.4.1 & Aplicar criterios para evaluar, seleccionar y hacer seguimiento a los proveedores externos. & Sin criterios documentados ni reevaluación tras los incumplimientos (acero entregado en 30~\%, Exp., p.~27; viguetas no disponibles, PAO2, p.~11). La certificación NB~997 e ISO~9001 de PRETENSA (p.~80) es un atributo del proveedor, no un criterio registrado. & 0 & --- \\
8.4.2 & Verificar lo suministrado antes de incorporarlo a la obra (certificados, inspección de recepción, ensayos). & Certificados de PRETENSA y UNICON del 27/11/2025, sin firma y posteriores al inicio del armado (14/11/2025) (PAO2, pp.~80--81, 122); certificados de acero de $\phi$20 a $\phi$8 (pp.~82--86), sin $\phi$6 ni $\phi$10; sin certificados de cemento ni áridos; sin inspección de recepción. & 0,5 & Cl.~26.6 (PAO2, p.~102); ACI~318-25 (p.~587). \\
8.4.3 & Comunicar a los proveedores los requisitos de lo que se compra (tipo, dimensiones, cantidad, aprobación). & No hay órdenes de compra; el EPS entregado (20~$\times$~100~$\times$~44~cm; 2.350 piezas) no coincide con la dimensión de la ET (PAO2, p.~81; ET, p.~3). & 0 & --- \\
8.5.1 a & Ejecutar en condiciones controladas: información disponible, seguimiento en las etapas apropiadas y personal competente. & Verificación del armado antes de cada vaciado, registrada en el Libro de Órdenes (PAO2, pp.~114--125), sin lista de verificación ni criterios; no consta la fecha del hormigonado de la losa (p.~125). & 0,5 & --- \\
8.5.1 b & Controlar el hormigón y validar el proceso de hormigonado: dosificación aprobada, asentamiento, muestreo, ensayos a 28 días y curado. & Dosificaciones no presentadas (PAO2, p.~9); instrucciones reiteradas sin registro (pp.~120, 121, 125); ninguna probeta de la losa; columnas y vigas ensayadas a 55--76 días (p.~78). & 0 & ACI~318-25 (pp.~577, 601--603). \\
8.5.2 & Identificar las salidas y su estado de inspección, y asegurar la trazabilidad de materiales y probetas. & Coladas identificadas en los certificados de acero (PAO2, pp.~82--86) sin vínculo con los elementos; probetas rotuladas solo «COLUMNA» o «VIGA» (p.~78). & 0 & --- \\
8.5.3 & Cuidar la propiedad del cliente o de proveedores externos. & --- & NA & La UATF no suministró materiales ni equipos para la losa. \\
8.5.4 & Preservar los materiales y el producto (acopio, protección, curado). & Plan de acopio limitado a «transportar el material y ubicarlo cerca de la obra» (PAO2, p.~10); material y restos dispersos en las fotografías (pp.~75--76); sin registro de curado. & 0 & Cl.~26.9 (PAO2, p.~102). \\
8.5.6 & Revisar y controlar los cambios no planificados en la ejecución. & Modificación estructural y cambio de complemento sin registro (PAO2, pp.~11, 81); vaciado «desde dos frentes» sin plan de juntas (p.~125). & 0 & --- \\
8.6 & Liberar cada etapa solo después de verificar los requisitos, con registro de quién autoriza. & Ocho autorizaciones de vaciado firmadas por la Supervisión (PAO2, pp.~114--125); la losa se certificó para pago el 08/12/2025 (p.~16) antes de que pudiera existir un ensayo a 28 días (26/12/2025). & 0,5 & --- \\
8.7 & Identificar y controlar las salidas no conformes y registrar su tratamiento. & Ningún rechazo, no conformidad ni retrabajo registrado en los 24 asientos (PAO2, pp.~114--125) ni en el informe de Supervisión (pp.~6--11). & 0 & --- \\
%
\multicolumn{5}{@{}l}{\textbf{9. Evaluación del desempeño}}\\
9.1.1 a & Determinar qué se mide (calidad, desperdicio, costos), con qué métodos y cuándo se analizan los resultados. & Solo se mide el avance físico-financiero por planillas (PAO2, pp.~10, 16). & 0,5 & Sin indicadores de calidad ni de desperdicio. \\
9.1.1 b & Medir y conciliar las cantidades ejecutadas (medición conjunta, neta y con respaldo) antes de cada planilla. & Cómputos firmados por contratista, Supervisión y Fiscal (PAO2, pp.~18--32), pero el acero de la losa sin planilla de armado (p.~20), correcciones posteriores (ítem~4: $-$5,66~\mtres{}; acero: $-$1.862,81~kg en la planilla de cierre firmada y $-$1.915,16~kg en su versión digital) y tres cifras distintas del acero acumulado (p.~16; PAO3; Exp., pp.~2, 12). & 0,5 & Cl.~27.ª (PAO2, p.~99). \\
9.1.2 & Hacer seguimiento de la percepción del cliente sobre el cumplimiento de sus requisitos. & El informe de Supervisión califica los trabajos de «buena calidad» y evalúa al Director de Obra (PAO2, pp.~9, 11), sin evidencia objetiva que respalde la calificación. & 0,5 & --- \\
9.1.3 & Analizar y evaluar los datos (conformidad, desempeño de proveedores, eficacia de las acciones). & Sin evidencia; el desfase del periodo se calculó mal: 3,03~\% en lugar de 1,07~\% (PAO2, p.~10). & 0 & --- \\
9.2 & Realizar auditorías internas planificadas. & No se encontraron. & 0 & --- \\
9.3 & Revisar el sistema por la dirección a intervalos planificados. & No se encontró. & 0 & --- \\
%
\multicolumn{5}{@{}l}{\textbf{10. Mejora}}\\
10.1 & Determinar y aplicar oportunidades de mejora. & La Supervisión recomendó presentar las garantías de los materiales, romper probetas y presentar dosificaciones (PAO2, p.~9); no consta que se atendieran en la losa (p.~78). & 0 & --- \\
10.2.1 a & Reaccionar ante las no conformidades: controlarlas, corregirlas y asumir sus consecuencias. & Se corrigieron cantidades después de pagadas: ítem~4, $-$5,66~\mtres{} (PAO2, pp.~16, 18); acero, $-$1.862,81~kg en la planilla de cierre (Exp., p.~2). & 0,5 & Correcciones de medición, sin análisis. \\
10.2.1 b & Analizar las causas, implementar acciones correctivas y verificar su eficacia. & Sin evidencia de análisis de causas ni de acciones correctivas. & 0 & --- \\
10.2.2 & Conservar información documentada de las no conformidades y de las acciones tomadas. & No hay registro de no conformidades. & 0 & --- \\
10.3 & Mejorar continuamente la conveniencia, adecuación y eficacia del sistema. & No se encontró. & 0 & --- \\
\end{longtable}
\vspace{-6pt}
\nota{Elaboración propia. Fuentes: PAO2 = \textcite{uatf2025pao2} (página del archivo PDF; incluye el contrato, el Libro de Órdenes y los certificados); PAO3 = \textcite{uatf_pao3} (archivo de planilla y cómputos de cierre); Exp. = \textcite{uatf_item13} (página del archivo PDF); ET = \textcite{uatf_et}; Plano~02 = sección y detalles de la losa (Anexo de planos). Calificación: 1 = cumple; 0,5 = cumple parcialmente; 0 = no cumple o sin evidencia; NA = no aplica. ACI~318-25 = \textcite{aci318_2025}.}
\endgroup

\textbf{Resumen por cláusula.} La Tabla~\ref{tab:lista-iso-resumen} resume la aplicación. Ningún requisito alcanzó la calificación de cumplimiento pleno: 19 se cumplieron parcialmente y 29 no se cumplieron, de un total de 48 requisitos aplicables.

\begin{table}[H]
\caption{Resumen de la lista de verificación C.1 por cláusula}
\label{tab:lista-iso-resumen}
\setlength{\tabcolsep}{4pt}
\begin{tabularx}{\linewidth}{@{}LQ{1.35cm}Q{1.0cm}Q{1.25cm}Q{1.25cm}Q{1.25cm}Q{1.3cm}Q{1.65cm}@{}}
\toprule
\textbf{Cláusula} & \textbf{Requisitos aplicables} & \textbf{NA} & \textbf{Cumple (1)} & \textbf{Parcial (0,5)} & \textbf{No cumple (0)} & \textbf{Puntaje} & \textbf{Cumplimiento} \\
\midrule
4. Contexto de la organización & 4 & 0 & 0 & 2 & 2 & 1,0 & 25,0~\% \\
5. Liderazgo & 5 & 0 & 0 & 2 & 3 & 1,0 & 20,0~\% \\
6. Planificación & 4 & 0 & 0 & 1 & 3 & 0,5 & 12,5~\% \\
7. Apoyo & 10 & 0 & 0 & 5 & 5 & 2,5 & 25,0~\% \\
8. Operación & 14 & 2 & 0 & 5 & 9 & 2,5 & 17,9~\% \\
9. Evaluación del desempeño & 6 & 0 & 0 & 3 & 3 & 1,5 & 25,0~\% \\
10. Mejora & 5 & 0 & 0 & 1 & 4 & 0,5 & 10,0~\% \\
\midrule
\textbf{Global} & \textbf{48} & \textbf{2} & \textbf{0} & \textbf{19} & \textbf{29} & \textbf{9,5} & \textbf{19,8~\%} \\
\bottomrule
\end{tabularx}
\nota{Elaboración propia. Cumplimiento = puntaje / requisitos aplicables $\times$ 100. El cumplimiento global se calcula con el total de requisitos (9,5/48); el promedio simple de las siete cláusulas es 19,3~\%. Según la escala de \textcite[p.~42]{mamanicalsina2025}, todas las cláusulas y el resultado global se ubican en el rango «deficiente o no implementado» (0--25~\%).}
\end{table}
}{}

% =====================================================================
\clearpage
\section*{C.2 Fichas PPI de liberación por paño}
\addcontentsline{toc}{section}{C.2 Fichas PPI de liberación por paño}

{\setstretch{1.15}\small
\textbf{Uso.} Una ficha por paño. Marque C (conforme), NC (no conforme) o NA (no aplica). Los puntos de espera (H) no se superan sin la firma del Supervisor; los de aviso (W) se notifican a la Supervisión. Toda NC se registra en el formato C.5 con su número. Criterios: Tabla~\ref{tab:ppi} del plan de la calidad.\par}

\cabeceraformato{C.2}{Ficha PPI de liberación del paño}

{\footnotesize\setstretch{1}\renewcommand{\arraystretch}{1.3}
\noindent\begin{tabularx}{\linewidth}{@{}|L|L|L|L|@{}}
\hline
\campo{Losa N.\textsuperscript{o}} & \campo{Paño N.\textsuperscript{o}} & \campo{Ejes} & \campo{Área (\mdos{})} \\
\campo{Viguetas (Plano 03): n.\textsuperscript{o}} & \campo{Tipo y longitud} & \campo{Luz libre (m)} & \campo{Fecha de inicio} \\
\hline
\end{tabularx}\par}

\newcommand{\ppiencabezado}[1]{%
\multicolumn{7}{|@{\,}l|}{\cellcolor{gray!15}\textbf{#1}} \\ \hline
\textbf{N.\textsuperscript{o}} & \textbf{Punto de control} & \textbf{Criterio de aceptación} & \textbf{C} & \textbf{NC} & \textbf{NA} & \textbf{Observación o valor medido} \\ \hline}
\newcommand{\ppifirmas}[2]{%
\multicolumn{7}{|@{\,}p{\dimexpr\linewidth-2\tabcolsep\relax}|}{\textbf{#1}\par #2\par\vspace{2pt}} \\ \hline}

{\footnotesize\setstretch{1}\renewcommand{\arraystretch}{1.25}\setlength{\tabcolsep}{3pt}
\begin{longtable}{@{}|P{0.75cm}|P{3.3cm}|P{5.2cm}|Q{0.45cm}|Q{0.55cm}|Q{0.55cm}|P{3.85cm}|@{}}
\hline
\ppiencabezado{A. Apuntalamiento y encofrado --- punto de espera (H)}
A1 & Ejes y niveles & Según el Plano~01; niveles de fondo de viga verificados & & & & \\ \hline
A2 & Encofrado de vigas & Estanco, arriostrado y limpio; alero $\geq 20$~cm donde la vigueta corre paralela a la viga & & & & \\ \hline
A3 & Soleras & Sección $\geq 2'' \times 4''$; una (luz de 2 a 4~m), dos (4,5 a 6~m) o tres (6,5 a 7,5~m), transversales a las viguetas & & & & N.\textsuperscript{o} de soleras: \\ \hline
A4 & Puntales & Diámetro $\geq 4''$; separación de 0,80 a 1,00~m; sobre tabla o cuña; arriostrados en cruz & & & & Separación medida (m): \\ \hline
A5 & Contraflecha & Según las instrucciones del fabricante & & & & Medida (mm): \\ \hline
\ppifirmas{Liberación A}{\firma{Ejecutó (Capataz)}\hfill\firma{Verificó (encargado de calidad)}\hfill\firma{Liberó (Supervisor)}\hfill\firma{Fecha y hora}}
\ppiencabezado{B. Viguetas y complementos EPS --- punto de aviso (W)}
B1 & Viguetas & Número, tipo y longitud según el Plano~03; lote aceptado (ficha C.3) & & & & C.3 N.\textsuperscript{o}: \\ \hline
B2 & Separación & 50~cm entre ejes & & & & \\ \hline
B3 & Apoyos & $\geq 5$~cm si se vacía con la viga; $\geq 10$~cm sobre viga ya vaciada & & & & Mínimo medido (cm): \\ \hline
B4 & Estado & Sin fisuras, golpes ni flecha & & & & \\ \hline
B5 & Complementos EPS & Geometría aprobada; piezas enteras y ajustadas; sin huecos; cortes solo en bordes y nervios & & & & Piezas colocadas: \\ \hline
B6 & Tránsito & Tablones o pasarelas colocados; no se pisan las piezas & & & & \\ \hline
\ppifirmas{Conformidad B}{\firma{Ejecutó (Capataz)}\hfill\firma{Verificó (encargado de calidad)}\hfill\firma{Fecha}}
\ppiencabezado{C. Nervios, armadura de distribución, negativos y continuidad --- punto de espera (H)}
C1 & Nervios transversales & Ubicación y dimensiones según el Plano~02 & & & & \\ \hline
C2 & Armadura de nervios & Según planos: número, diámetro y longitud & & & & \\ \hline
C3 & Parrilla & $\phi$6~c/20~cm en ambas direcciones (Plano~02) & & & & Separación medida (cm): \\ \hline
C4 & Negativos y continuidad & Según planos y lo exigido por el fabricante de viguetas & & & & \\ \hline
C5 & Separadores y recubrimiento & Parrilla separada del EPS; separadores suficientes; recubrimiento según planos & & & & Recubrimiento (cm): \\ \hline
C6 & Amarres y limpieza & Empalmes y amarres completos; barras sin barro, aceite ni óxido suelto & & & & \\ \hline
C7 & Trazabilidad del acero & Certificado por colada y diámetro aceptado (ficha C.3) & & & & C.3 N.\textsuperscript{o}: \\ \hline
\ppifirmas{Liberación C}{\firma{Ejecutó (Capataz)}\hfill\firma{Verificó (encargado de calidad)}\hfill\firma{Liberó (Supervisor)}\hfill\firma{Fecha y hora}}
\ppiencabezado{D. Instalaciones embebidas --- punto de aviso (W)}
D1 & Eléctricas & Ductos y cajas según planos de instalaciones & & & & \\ \hline
D2 & Sanitarias y pases & Según planos; pases sellados & & & & \\ \hline
D3 & Integridad & Sin cortes en viguetas ni en nervios & & & & \\ \hline
\ppifirmas{Conformidad D}{\firma{Instalador}\hfill\firma{Verificó (encargado de calidad)}\hfill\firma{Fecha}}
\ppiencabezado{E. Liberación para el vaciado y control del vaciado --- punto de espera (H)}
E1 & Etapas previas & A, B, C y D conformes; ninguna NC abierta en el paño & & & & \\ \hline
E2 & Superficie & Limpia, sin escombros, hielo ni agua estancada; humedecida & & & & \\ \hline
E3 & Dosificación y existencias & Dosificación aprobada; cemento y áridos para el vaciado completo & & & & Bolsas disponibles / requeridas: \\ \hline
E4 & Equipos y personal & Mezcladora, vibrador y vibrador de reserva; técnico de muestreo y moldes & & & & N.\textsuperscript{o} de muestras previstas: \\ \hline
E5 & Clima & Pronóstico sin lluvia; si la temperatura mínima prevista es menor de 4,4~$^{\circ}$C, plan de tiempo frío activo & & & & T mín. prevista ($^{\circ}$C): \\ \hline
E6 & Asentamiento & Dentro del rango de la dosificación (ASTM~C143) & & & & Valores (cm): \\ \hline
E7 & Temperaturas & Hormigón registrada al muestrear; agua de amasado $> 5$~$^{\circ}$C & & & & T hormigón / agua ($^{\circ}$C): \\ \hline
E8 & Colocación & Por franjas, continua, sin juntas frías; vibrado; tránsito sobre pasarelas sin pisar la parrilla & & & & \\ \hline
E9 & Registro del vaciado & Hora de inicio y fin; volumen; pastones; bolsas de cemento; fotografía fechada de la toma de muestras & & & & Inicio / fin / \mtres{} / pastones / bolsas: \\ \hline
\ppifirmas{Liberación E (antes de vaciar)}{\firma{Encargado de calidad}\hfill\firma{Liberó (Supervisor)}\hfill\firma{Libro de Órdenes, hoja N.\textsuperscript{o}}\hfill\firma{Fecha y hora}}
\end{longtable}}

\noindent{\small\setstretch{1.1}\textbf{F. Registro de curado (W).} Curado mínimo de 7 días con el hormigón a 10~$^{\circ}$C o más; si no se cumple la temperatura, se prolonga el curado.\par}
\smallskip
{\footnotesize\setstretch{1}\renewcommand{\arraystretch}{1.3}
\noindent\begin{tabularx}{\linewidth}{@{}|Q{1.0cm}|P{2.0cm}|L|P{2.0cm}|P{2.0cm}|P{2.2cm}|P{2.0cm}|@{}}
\hline
\textbf{Día} & \textbf{Fecha} & \textbf{Método (riego, arpillera, plástico)} & \textbf{T mín. ambiente ($^{\circ}$C)} & \textbf{Protección contra helada} & \textbf{Responsable} & \textbf{Foto N.\textsuperscript{o}} \\ \hline
1 & & & & & & \\ \hline
2 & & & & & & \\ \hline
3 & & & & & & \\ \hline
4 & & & & & & \\ \hline
5 & & & & & & \\ \hline
6 & & & & & & \\ \hline
7 & & & & & & \\ \hline
\fila & & & & & & \\ \hline
\end{tabularx}\par}

\medskip
\noindent{\small\setstretch{1.1}\textbf{G. Desapuntalamiento (H) y aceptación del paño.}\par}
\smallskip
{\footnotesize\setstretch{1}\renewcommand{\arraystretch}{1.3}
\noindent\begin{tabularx}{\linewidth}{@{}|P{0.75cm}|P{4.0cm}|L|Q{0.45cm}|Q{0.55cm}|P{3.4cm}|@{}}
\hline
\textbf{N.\textsuperscript{o}} & \textbf{Punto de control} & \textbf{Criterio de aceptación} & \textbf{C} & \textbf{NC} & \textbf{Valor} \\ \hline
G1 & Plazo desde el vaciado & $\geq 15$ días (luz $< 5$~m) o $\geq 21$ días (luz $> 5$~m) & & & Días: \\ \hline
G2 & Cilindros curados en obra & Si la Supervisión los exige: todos $\geq$ la resistencia requerida para la etapa & & & kg/cm\textsuperscript{2}: \\ \hline
G3 & Cargas y estado & Sin cargas de construcción excesivas; sin fisuras ni flechas visibles después del retiro & & & \\ \hline
G4 & Aceptación para pago & Ensayos de aceptación del vaciado conformes; conciliación C.4 cerrada; ninguna NC abierta & & & Ensayos N.\textsuperscript{o}: \\ \hline
\multicolumn{6}{|@{\,}p{\dimexpr\linewidth-2\tabcolsep\relax}|}{\textbf{Liberación G}\par\firma{Encargado de calidad}\hfill\firma{Liberó (Supervisor)}\hfill\firma{Fecha}\par\vspace{2pt}} \\ \hline
\end{tabularx}\par}

% =====================================================================
\clearpage
\section*{C.3 Ficha de recepción de materiales}
\addcontentsline{toc}{section}{C.3 Ficha de recepción de materiales}

{\setstretch{1.15}\small
\textbf{Uso.} Una ficha por entrega y por material, llenada \emph{antes} de colocar el material (PCM-01). Ningún material se coloca sin esta ficha aceptada. Complete solo el bloque del material recibido. Criterios: Tabla~\ref{tab:pcm01}.\par}

\cabeceraformato{C.3}{Ficha de recepción de materiales}

{\footnotesize\setstretch{1}\renewcommand{\arraystretch}{1.3}
\noindent\begin{tabularx}{\linewidth}{@{}|L|L|L|@{}}
\hline
\campo{Ficha N.\textsuperscript{o}} & \campo{Fecha y hora} & \campo{Proveedor} \\
\campo{Material} & \campo{Nota de remisión N.\textsuperscript{o}} & \campo{Factura N.\textsuperscript{o}} \\
\campo{Orden de compra N.\textsuperscript{o}} & \campo{Lote o colada} & \campo{Destino (losa y paños)} \\
\campo{Cantidad remitida} & \campo{Cantidad recibida (conteo)} & \campo{Unidad} \\
\hline
\end{tabularx}\par}

\medskip
\noindent{\small\textbf{1. Verificación documental}}\par\smallskip
{\footnotesize\setstretch{1}\renewcommand{\arraystretch}{1.3}
\noindent\begin{tabularx}{\linewidth}{@{}|L|Q{1.0cm}|Q{1.0cm}|Q{1.0cm}|P{4.2cm}|@{}}
\hline
\textbf{Documento} & \textbf{Sí} & \textbf{No} & \textbf{NA} & \textbf{Observación} \\ \hline
Certificado de calidad firmado, con lote y cantidad & & & & \\ \hline
Fecha del certificado anterior a la colocación & & & & \\ \hline
Viguetas: conformidad con la NB~997 declarada en el certificado & & & & \\ \hline
Acero: certificado por colada para cada diámetro entregado ($\phi$6, $\phi$8, $\phi$10, $\phi$12\ldots) & & & & \\ \hline
Cemento: sello IBNORCA, tipo, clase y fecha de fabricación & & & & \\ \hline
Remisión, factura y orden de compra coinciden en material y cantidad & & & & \\ \hline
\end{tabularx}\par}

\medskip
\noindent{\small\textbf{2. Inspección física por material}}\par\smallskip
{\footnotesize\setstretch{1}\renewcommand{\arraystretch}{1.3}\setlength{\tabcolsep}{3pt}
\noindent\textbf{a) Viguetas pretensadas} (muestra: cinco viguetas por entrega y tipo; inspección visual de todas)\par\smallskip
\noindent\begin{tabularx}{\linewidth}{@{}|Q{0.6cm}|Q{1.3cm}|Q{1.5cm}|Q{1.5cm}|Q{1.4cm}|Q{2.2cm}|Q{2.0cm}|L|Q{0.9cm}|@{}}
\hline
\textbf{N.\textsuperscript{o}} & \textbf{Tipo} & \textbf{L teórica (mm)} & \textbf{L medida (mm)} & \textbf{$\Delta L$ (+30/$-$20)} & \textbf{Sección $a$, $b$, $h$ (+5/$-$2~mm)} & \textbf{Contraflecha y flecha lateral ($\leq L/300$)} & \textbf{Fisuras, oquedades, armadura visible} & \textbf{C / NC} \\ \hline
1 & & & & & & & & \\ \hline
2 & & & & & & & & \\ \hline
3 & & & & & & & & \\ \hline
4 & & & & & & & & \\ \hline
5 & & & & & & & & \\ \hline
\multicolumn{9}{|@{\,}l|}{\fila Defectos principales: \underline{\hspace{1.2cm}} (rechazo del lote si $> 2$) \qquad Defectos secundarios: \underline{\hspace{1.2cm}} (rechazo del lote si $> 4$)} \\ \hline
\end{tabularx}

\medskip
\noindent\textbf{b) Complementos EPS} (muestra: cinco piezas; densidad = masa / volumen; declarada: 10~kg/\mtres{})\par\smallskip
\noindent\begin{tabularx}{\linewidth}{@{}|Q{0.6cm}|Q{1.6cm}|Q{1.6cm}|Q{1.6cm}|Q{1.6cm}|Q{1.8cm}|Q{1.9cm}|L|Q{0.9cm}|@{}}
\hline
\textbf{N.\textsuperscript{o}} & \textbf{Largo (cm)} & \textbf{Ancho (cm)} & \textbf{Alto (cm)} & \textbf{Masa (kg)} & \textbf{Volumen (\mtres{})} & \textbf{Densidad (kg/\mtres{})} & \textbf{Estado (roturas, deformación)} & \textbf{C / NC} \\ \hline
1 & & & & & & & & \\ \hline
2 & & & & & & & & \\ \hline
3 & & & & & & & & \\ \hline
4 & & & & & & & & \\ \hline
5 & & & & & & & & \\ \hline
\end{tabularx}

\medskip
\noindent\textbf{c) Acero de refuerzo} (pesaje de tres barras por diámetro y colada; masa medida $\geq 94$~\% de la nominal)\par\smallskip
\noindent\begin{tabularx}{\linewidth}{@{}|Q{0.9cm}|Q{1.6cm}|Q{1.3cm}|Q{1.7cm}|Q{1.5cm}|Q{1.5cm}|Q{1.6cm}|Q{1.3cm}|L|Q{0.9cm}|@{}}
\hline
\textbf{$\phi$ (mm)} & \textbf{Colada} & \textbf{N.\textsuperscript{o} de barras} & \textbf{Masa de 3 barras (kg)} & \textbf{Longitud total (m)} & \textbf{Masa por metro (kg/m)} & \textbf{Nominal (kg/m)} & \textbf{\% de la nominal} & \textbf{Corrugado, óxido, aceite} & \textbf{C / NC} \\ \hline
\fila & & & & & & & & & \\ \hline
\fila & & & & & & & & & \\ \hline
\fila & & & & & & & & & \\ \hline
\fila & & & & & & & & & \\ \hline
\end{tabularx}

\medskip
\noindent\textbf{d) Cemento, áridos y agua}\par\smallskip
\noindent\begin{tabularx}{\linewidth}{@{}|P{2.2cm}|L|L|Q{0.9cm}|@{}}
\hline
\textbf{Material} & \textbf{Verificación} & \textbf{Resultado} & \textbf{C / NC} \\ \hline
Cemento & Bolsas recibidas; bolsas rotas o con grumos, terrones o costras; fecha de fabricación & & \\ \hline
Arena y grava & Volquete N.\textsuperscript{o}; volumen cubicado (\mtres{}); limpieza (sin arcilla ni materia orgánica); tamaño máximo $\leq$ mitad del espesor de la carpeta & & \\ \hline
Agua & Fuente aprobada; limpieza; temperatura al amasar $> 5$~$^{\circ}$C & & \\ \hline
\end{tabularx}\par}

\medskip
\noindent{\small\textbf{3. Decisión}}\par\smallskip
{\footnotesize\setstretch{1}\renewcommand{\arraystretch}{1.4}
\noindent\begin{tabularx}{\linewidth}{@{}|L|@{}}
\hline
\casilla\ Aceptado \qquad \casilla\ Aceptado bajo concesión (autorización escrita de la Supervisión N.\textsuperscript{o} \underline{\hspace{1.8cm}}) \qquad \casilla\ Rechazado $\rightarrow$ NC N.\textsuperscript{o} \underline{\hspace{1.5cm}} (formato C.5) \\
Identificación del lote aceptado (tarjeta verde) con el código \underline{\hspace{2.5cm}} \qquad Ingreso al kardex: fecha \underline{\hspace{2.2cm}} \\
\firma{Recibió (Almacenero)}\hfill\firma{Inspeccionó y decidió (encargado de calidad)}\hfill\firma{Presenció (Supervisor)}\hfill\firma{Fecha} \\[2pt]
\hline
\end{tabularx}\par}

% =====================================================================
\clearpage
\section*{C.4 Ficha de control de materiales y conciliación por paño}
\addcontentsline{toc}{section}{C.4 Ficha de control de materiales y conciliación por paño}

{\setstretch{1.15}\small
\textbf{Uso.} La parte (a) se emite por cada despacho o devolución; la parte (b) se lleva por material en el almacén; la parte (c) se cierra al terminar cada paño y la parte (d) antes de cada planilla de avance (PCM-03 y PCM-04). Fórmulas: $Q_C = Q_D - Q_V$; $IDM = (Q_C - Q_T)/Q_T \times 100$; $\delta_{X/Y} = (Q_X - Q_Y)/Q_Y \times 100$; umbrales de alerta de la Tabla~\ref{tab:tolerancias}.\par}

\cabeceraformato{C.4}{Control de materiales y conciliación por paño}

\noindent{\small\textbf{(a) Vale de salida o de devolución}}\par\smallskip
{\footnotesize\setstretch{1}\renewcommand{\arraystretch}{1.3}\setlength{\tabcolsep}{3pt}
\noindent\begin{tabularx}{\linewidth}{@{}|Q{1.1cm}|Q{1.6cm}|Q{1.2cm}|L|Q{1.2cm}|Q{1.5cm}|Q{1.7cm}|P{2.0cm}|P{2.0cm}|@{}}
\hline
\textbf{Vale N.\textsuperscript{o}} & \textbf{Fecha} & \textbf{Paño} & \textbf{Material (tipo, $\phi$, longitud)} & \textbf{Unidad} & \textbf{Cantidad} & \textbf{Salida o devolución} & \textbf{Entregó (firma)} & \textbf{Recibió (firma)} \\ \hline
\fila & & & & & & & & \\ \hline
\fila & & & & & & & & \\ \hline
\fila & & & & & & & & \\ \hline
\fila & & & & & & & & \\ \hline
\fila & & & & & & & & \\ \hline
\end{tabularx}\par}

\medskip
\noindent{\small\textbf{(b) Kardex del material:} \underline{\hspace{4cm}} \quad Unidad: \underline{\hspace{2cm}}}\par\smallskip
{\footnotesize\setstretch{1}\renewcommand{\arraystretch}{1.3}\setlength{\tabcolsep}{3pt}
\noindent\begin{tabularx}{\linewidth}{@{}|Q{1.6cm}|L|Q{1.5cm}|Q{1.5cm}|Q{1.6cm}|Q{1.5cm}|Q{1.8cm}|Q{1.6cm}|@{}}
\hline
\textbf{Fecha} & \textbf{Documento (C.3, vale, devolución)} & \textbf{Entrada} & \textbf{Salida} & \textbf{Devolución} & \textbf{Saldo} & \textbf{Inventario físico} & \textbf{Diferencia} \\ \hline
\fila & & & & & & & \\ \hline
\fila & & & & & & & \\ \hline
\fila & & & & & & & \\ \hline
\fila & & & & & & & \\ \hline
\fila & & & & & & & \\ \hline
\fila & & & & & & & \\ \hline
\end{tabularx}\par}

\medskip
\noindent{\small\textbf{(c) Conciliación del paño} \quad Losa N.\textsuperscript{o} \underline{\hspace{1cm}} \quad Paño N.\textsuperscript{o} \underline{\hspace{1cm}} \quad Área (\mdos{}) \underline{\hspace{1.4cm}} \quad Fechas \underline{\hspace{3cm}}}\par\smallskip
{\footnotesize\setstretch{1}\renewcommand{\arraystretch}{1.3}\setlength{\tabcolsep}{2.5pt}
\noindent\begin{tabularx}{\linewidth}{@{}|L|Q{0.9cm}|Q{1.3cm}|Q{1.3cm}|Q{1.3cm}|Q{1.4cm}|Q{1.2cm}|Q{0.9cm}|Q{1.3cm}|Q{1.2cm}|Q{0.9cm}|@{}}
\hline
\textbf{Material} & \textbf{Unidad} & \textbf{$Q_T$ teórica (planos)} & \textbf{$Q_D$ despachada} & \textbf{$Q_V$ devuelta} & \textbf{$Q_C = Q_D - Q_V$} & \textbf{$IDM$ (\%)} & \textbf{Meta (\%)} & \textbf{$Q_M$ medida} & \textbf{$\delta_{M/T}$ (\%)} & \textbf{Alerta (S/N)} \\ \hline
Viguetas tipo \underline{\hspace{0.6cm}} & m & & & & & & 1 & & & \\ \hline
Viguetas tipo \underline{\hspace{0.6cm}} & m & & & & & & 1 & & & \\ \hline
Complemento EPS & pza & & & & & & 3 & & & \\ \hline
Acero $\phi$6 (parrilla) & kg & & & & & & 4 & & & \\ \hline
Acero $\phi$ \underline{\hspace{0.5cm}} (nervios) & kg & & & & & & 4 & & & \\ \hline
Acero $\phi$ \underline{\hspace{0.5cm}} (negativos) & kg & & & & & & 4 & & & \\ \hline
Hormigón H25 & \mtres{} & & & & & & 3 & & & \\ \hline
Cemento & bolsa & & & & & & & & & \\ \hline
Arena & \mtres{} & & & & & & & & & \\ \hline
Grava & \mtres{} & & & & & & & & & \\ \hline
Alambre de amarre & kg & & & & & & & & & \\ \hline
\end{tabularx}\par}

\medskip
\noindent{\small\textbf{(d) Conciliación acumulada para la planilla de avance N.\textsuperscript{o}} \underline{\hspace{1cm}} \quad Periodo \underline{\hspace{3.5cm}}}\par\smallskip
{\footnotesize\setstretch{1}\renewcommand{\arraystretch}{1.3}\setlength{\tabcolsep}{2.5pt}
\noindent\begin{tabularx}{\linewidth}{@{}|L|Q{1.25cm}|Q{1.25cm}|Q{1.25cm}|Q{1.1cm}|Q{1.1cm}|Q{1.25cm}|Q{1.25cm}|Q{1.1cm}|Q{1.0cm}|Q{0.9cm}|@{}}
\hline
\textbf{Ítem o material (ítem de pago)} & \textbf{$Q_T$ acum.} & \textbf{$Q_E$ acum. (remisiones)} & \textbf{$Q_C$ acum. (vales)} & \textbf{$I_F$ físico} & \textbf{$\Delta_I$} & \textbf{$Q_M$ acum. (medición)} & \textbf{$Q_P$ acum. (planilla)} & \textbf{$\delta_{P/M}$ (\%)} & \textbf{kg/\mdos{}} & \textbf{Alerta} \\ \hline
Ítem 13 (\mdos{}) & & --- & --- & --- & --- & & & & --- & \\ \hline
Acero de la losa (kg); ítem \underline{\hspace{0.5cm}} & & & & & & & & & & \\ \hline
Viguetas (m) & & & & & & --- & --- & --- & --- & \\ \hline
Complemento EPS (pza) & & & & & & --- & --- & --- & --- & \\ \hline
Cemento (bolsas) & & & & & & --- & --- & --- & --- & \\ \hline
Hormigón (\mtres{}) & & & & --- & --- & & & & --- & \\ \hline
\multicolumn{11}{|@{\,}p{\dimexpr\linewidth-2\tabcolsep\relax}|}{\fila Coherencia entre documentos (misma cifra en planilla, cómputo e informes): \casilla\ Sí \casilla\ No \qquad Acero de la losa pagado en un solo ítem: \casilla\ Sí \casilla\ No} \\ \hline
\multicolumn{11}{|@{\,}p{\dimexpr\linewidth-2\tabcolsep\relax}|}{\firma{Elaboró (encargado de calidad)}\hfill\firma{Revisó (Almacenero)}\hfill\firma{Aprobó (Supervisor)}\hfill\firma{Fecha}\par\vspace{2pt}} \\ \hline
\end{tabularx}\par}

% =====================================================================
\clearpage
\section*{C.5 Registro de no conformidades y acción correctiva}
\addcontentsline{toc}{section}{C.5 Registro de no conformidades y acción correctiva}

{\setstretch{1.15}\small
\textbf{Uso.} Un registro por no conformidad (NC), abierto dentro de las 24~h de su detección y cerrado en 7 días o menos (objetivo OC9). Procedimiento: Figura~\ref{fig:flujo-nc}. El costo se traslada a la planilla C.6.\par}

\cabeceraformato{C.5}{Registro de no conformidad y acción correctiva}

{\footnotesize\setstretch{1}\renewcommand{\arraystretch}{1.35}
\noindent\begin{tabularx}{\linewidth}{@{}|L|L|L|@{}}
\hline
\campo{NC N.\textsuperscript{o}} & \campo{Fecha de detección} & \campo{Detectada por} \\
\campo{Losa / paño / lote} & \campo{Requisito incumplido (plano, ET, PPI, PCM, contrato)} & \campo{Referencia} \\
\multicolumn{3}{|@{\,}l|}{Proceso: \casilla\ Recepción \casilla\ Almacén \casilla\ Ejecución del paño \casilla\ Hormigón \casilla\ Cantidades \casilla\ Documentos \quad Clasificación: \casilla\ Crítica \casilla\ Mayor \casilla\ Menor} \\
\hline
\end{tabularx}

\medskip
\noindent\begin{tabularx}{\linewidth}{@{}|L|@{}}
\hline
\textbf{1. Descripción de la no conformidad} (qué se encontró, dónde, magnitud; fotografía N.\textsuperscript{o}) \\
\fila \\ \fila \\ \hline
\textbf{2. Contención y corrección} \quad Paño o lote retenido: \casilla\ Sí \casilla\ No \quad Acción inmediata, responsable y fecha: \\
\fila \\ \hline
\textbf{3. Análisis de causa.} \emph{5 porqués:} \\
1.~¿Por qué? \hrulefill \\
2.~¿Por qué? \hrulefill \\
3.~¿Por qué? \hrulefill \\
4.~¿Por qué? \hrulefill \\
5.~¿Por qué? \hrulefill \\
\emph{Ishikawa (6M):} \casilla\ Mano de obra \casilla\ Método \casilla\ Materiales \casilla\ Maquinaria \casilla\ Medición \casilla\ Medio ambiente \qquad Causa raíz: \hrulefill \\
\hline
\end{tabularx}

\medskip
\noindent\begin{tabularx}{\linewidth}{@{}|Q{0.7cm}|L|P{3.0cm}|P{2.2cm}|P{3.0cm}|@{}}
\hline
\multicolumn{5}{|@{\,}l|}{\textbf{4. Acción correctiva} (elimina la causa; la aprueba el Director de obra)} \\ \hline
\textbf{N.\textsuperscript{o}} & \textbf{Acción} & \textbf{Responsable} & \textbf{Plazo} & \textbf{Recursos} \\ \hline
1 & & & & \\ \hline
2 & & & & \\ \hline
3 & & & & \\ \hline
\end{tabularx}

\medskip
\noindent\begin{tabularx}{\linewidth}{@{}|L|L|@{}}
\hline
\multicolumn{2}{|@{\,}l|}{\textbf{5. Verificación de la eficacia}} \\ \hline
\campo{Fecha de verificación} & \campo{Evidencia (paños o lotes revisados)} \\
\multicolumn{2}{|@{\,}l|}{¿Se repitió la NC? \casilla\ No (eficaz) \quad \casilla\ Sí (no eficaz: nuevo análisis de causa, punto 3)} \\ \hline
\multicolumn{2}{|@{\,}l|}{\textbf{6. Costo asociado (precios constantes del contrato)} \quad Categoría: \casilla\ Falla interna \casilla\ Falla externa} \\ \hline
\campo{Mano de obra: horas-hombre $\times$ costo horario = Bs} & \campo{Materiales: cantidad $\times$ precio = Bs} \\
\campo{Equipo y ensayos = Bs} & \campo{\textbf{Total de la NC = Bs}} \\ \hline
\multicolumn{2}{|@{\,}l|}{\textbf{7. Cierre} \quad Fecha de cierre \underline{\hspace{2.2cm}} \quad Días desde la detección \underline{\hspace{1.2cm}} (meta: $\leq 7$) \quad Actualizar: \casilla\ Riesgos \casilla\ PPI \casilla\ Procedimiento} \\
\multicolumn{2}{|@{\,}p{\dimexpr\linewidth-2\tabcolsep\relax}|}{\firma{Encargado de calidad}\hfill\firma{Director de obra}\hfill\firma{Supervisor (si es crítica o mayor)}\hfill\firma{Fecha}\par\vspace{2pt}} \\ \hline
\end{tabularx}\par}

\medskip
\noindent{\small\textbf{Registro general de no conformidades de la losa N.\textsuperscript{o}} \underline{\hspace{1cm}}}\par\smallskip
{\footnotesize\setstretch{1}\renewcommand{\arraystretch}{1.3}\setlength{\tabcolsep}{3pt}
\noindent\begin{tabularx}{\linewidth}{@{}|Q{0.8cm}|Q{1.5cm}|Q{1.3cm}|Q{1.3cm}|L|Q{1.8cm}|Q{1.5cm}|Q{0.9cm}|Q{1.4cm}|@{}}
\hline
\textbf{NC N.\textsuperscript{o}} & \textbf{Fecha} & \textbf{Paño o lote} & \textbf{Clase} & \textbf{Descripción breve} & \textbf{Responsable} & \textbf{Fecha de cierre} & \textbf{Días} & \textbf{Costo (Bs)} \\ \hline
\fila & & & & & & & & \\ \hline
\fila & & & & & & & & \\ \hline
\fila & & & & & & & & \\ \hline
\fila & & & & & & & & \\ \hline
\fila & & & & & & & & \\ \hline
\end{tabularx}\par}

